\exercisesection{}

\begin{enumerate}
\def\labelenumi{\arabic{enumi})}
\tightlist
\item
  Let \(\mathbf{E}\) be a set, \(\Phi\) a clan of subsets of \(\mathbf{E}\) , and \(m\) an additive set function defined on \(\Phi\) . Denote by \(\mathcal{E}(\Phi)\) the vector space of real \(\Phi\) -step functions on \(\mathbf{E}\) , and by \(J\) the linear form on \(\mathcal{E}(\Phi)\) such that \(J(\varphi_{\mathbf{A}}) = m(\mathbf{A})\) for all \(A \in \Phi\) .
\end{enumerate}

\begin{enumerate}
\def\labelenumi{\alph{enumi})}
\tightlist
\item
  For every \(\mathbf{A} \in \Phi\) , one sets
\end{enumerate}

\[
m _ {+} (\mathrm{A}) = \sup _ {\mathrm{B} \in \Phi , \mathrm{B} \subset \mathrm{A}} m (\mathrm{B}) \quad \text { and } \quad m _ {-} (\mathrm{A}) = \sup _ {\mathrm{B} \in \Phi , \mathrm{B} \subset \mathrm{A}} (- m (\mathrm{B})).
\]

Show that \(|m|(\mathbf{A}) = m_{+}(\mathbf{A}) + m_{-}(\mathbf{A})\) is the supremum of the set of numbers of the form \(\sum_{i=1}^{p}|m(\mathbf{A}_i)|\) for all the finite partitions \((\mathbf{A}_i)_{1\leqslant i\leqslant p}\) of A into sets belonging to \(\Phi\) .

\begin{enumerate}
\def\labelenumi{\alph{enumi})}
\setcounter{enumi}{1}
\item
  One says that \(m\) is relatively bounded if the set of numbers of the form \(m(\mathbf{B})\) with \(\mathbf{B} \in \Phi\) , \(\mathbf{B} \subset \mathbf{A}\) is bounded for any set \(\mathbf{A} \in \Phi\) . Show that the following conditions are equivalent: \(\alpha)\) \(m\) is relatively bounded; \(\beta)\) the linear form \(J\) on the Riesz space \(\mathcal{E}(\Phi)\) is relatively bounded; \(\gamma)\) \(m\) is the difference of two positive additive set functions on \(\Phi\) ; \(\delta)\) \(|m|(\mathbf{A})\) is finite for all \(\mathbf{A} \in \Phi\) (argue in a circle \(\alpha) \Rightarrow \beta) \Rightarrow \gamma) \Rightarrow \delta) \Rightarrow \alpha)\)).
\item
  Suppose that \(m\) is relatively bounded. Show that the set functions \(|m|\) , \(m_{+}\) and \(m_{-}\) are additive and that \(m = m_{+} - m_{-}\) (make use of Th. 1 of Ch. II, §2, No.~2). Moreover, if \(m'\) and \(m''\) are additive and positive set functions on \(\Phi\) such that \(m = m' - m''\) , then \(m' \geqslant m_{+}\) and \(m'' \geqslant m_{-}\) .
\end{enumerate}

\(d)\) One says that \(m\) is countably additive if \(m(A) = \sum_{n\geqslant 1}m(A_n)\) for every countable

partition \((\mathbf{A}_n)_{n\geqslant 1}\) of a set \(\mathbf{A}\in \Phi\) into sets belonging to \(\Phi\) . Show that this condition signifies that \(\lim_{n\to \infty}m(\mathbf{A}_n) = 0\) for every decreasing sequence \((\mathbf{A}_n)_{n\geqslant 1}\) of elements of \(\Phi\) such that \(\bigcap_{n\geqslant 1}\mathbf{A}_n = \emptyset\) . If \(m\) is countably additive, show that \(|m|\) is countably additive.

\begin{enumerate}
\def\labelenumi{\arabic{enumi})}
\setcounter{enumi}{1}
\tightlist
\item
  Let E be a set, \(\Phi\) a clan of subsets of E, and V a fully lattice-ordered Riesz space. A mapping \(m\) of \(\Phi\) into V is said to be positive if \(m(A) \geqslant 0\) for all \(A \in \Phi\) , that it is additive if \(m(A \cup B) = m(A) + m(B)\) when the sets \(A \in \Phi\) and \(B \in \Phi\) are disjoint, and that it is relatively bounded if the set of elements \(m(B)\) for \(B \in \Phi\) , \(B \subset A\) is bounded above in V for any \(A \in \Phi\) .
\end{enumerate}

\begin{enumerate}
\def\labelenumi{\alph{enumi})}
\tightlist
\item
  Let \(m\) be a relatively bounded mapping of \(\Phi\) into \(V\) . One defines for every \(A\) the elements \(m_{+}(A), m_{-}(A)\) and \(|m|(A)\) of \(V\) as in Exer. 1 a). Show that \(|m|(A)\)
\end{enumerate}

is the supremum of the set of elements of V of the form \(\sum_{i=1}^{p}|m(A_i)|\) for all the finite partitions \((\mathbf{A}_i)_{1\leqslant i\leqslant p}\) of A into sets belonging to \(\Phi\) . From this, deduce that the mapping \(|m|\) of \(\Phi\) into V is additive, then that \(m\) is the difference of the additive positive mappings \(m_{+}\) and \(m_{-}\) of \(\Phi\) into V.

\begin{enumerate}
\def\labelenumi{\alph{enumi})}
\setcounter{enumi}{1}
\tightlist
\item
  Let \(m'\) and \(m''\) be two additive positive mappings of \(\Phi\) into \(V\) and let \(m = m' - m''\) . Show that \(m\) is relatively bounded, and that \(m' \geqslant m_+\) and \(m'' \geqslant m_-\) . c) Generalize Exer. 1 d).
\end{enumerate}

\begin{enumerate}
\def\labelenumi{\arabic{enumi})}
\setcounter{enumi}{2}
\tightlist
\item
  Let \(\mathbf{E}\) be a set and \(\mathfrak{T}\) a tribe of subsets of \(\mathbf{E}\) . Denote by \(\mathcal{M}\) the set of mappings \(f\) of \(\mathbf{E}\) into \(\overline{\mathbf{R}}_+\) such that the set of \(x \in \mathbf{E}\) for which \(f(x) \geqslant c\) belongs to \(\mathfrak{T}\) for every \(c \in \overline{\mathbf{R}}_+\) .
\end{enumerate}

\begin{enumerate}
\def\labelenumi{\alph{enumi})}
\item
  Show that the limit superior and limit inferior of every sequence of elements of \(\mathcal{M}\) belong to \(\mathcal{M}\) (first study sequences that are increasing or decreasing).
\item
  Show that \(\mathfrak{T}\) is the set of subsets of \(\mathbf{E}\) whose characteristic function belongs to \(\mathcal{M}\) .
\item
  Let \(f \in \mathcal{M}\) . For every Borel subset \(B\) of \(\overline{\mathbf{R}}_+\) , the set \(f^{-1}(B)\) belongs to \(\mathfrak{T}\) (the set of subsets \(B\) of \(\overline{\mathbf{R}}_+\) such that \(f^{-1}(B) \in \mathfrak{T}\) is a tribe, containing the intervals \([c, +\infty]\) ).
\item
  Let \(f_1, \ldots, f_n\) be in \(\mathcal{M}\) and let \(\varphi\) be a Borel mapping of \((\overline{\mathbf{R}}_+)^n\) into \(\overline{\mathbf{R}}_+\) ; show that the mapping \(x \mapsto \varphi(f_1(x), \ldots, f_n(x))\) of E into \(\overline{\mathbf{R}}_+\) belongs to \(\mathcal{M}\) . Deduce from this that \(\mathcal{M}\) contains the sum of every series of elements of \(\mathcal{M}\) .
\item
  Show that \(\mathcal{M}\) is the set of limits of increasing sequences of finite positive \(\mathfrak{T}\) -step functions.
\end{enumerate}

\begin{enumerate}
\def\labelenumi{\arabic{enumi})}
\setcounter{enumi}{3}
\tightlist
\item
  Notations and hypotheses are those of the preceding exercise. One calls abstract measure on \((\mathbf{E},\mathfrak{T})\) every mapping \(m\) of \(\mathfrak{T}\) into \(\overline{\mathbf{R}}_+\) having the following property: for every sequence \((\mathrm{A}_n)_{n\in \mathbf{N}}\) of sets belonging to \(\mathfrak{T}\) , pairwise disjoint, one has \(m\big(\bigcup_{n\in \mathbf{N}}\mathrm{A}_n\big) =\)
\end{enumerate}

\(\sum_{n\in \mathbf{N}}m(\mathrm{A}_n)\) . One calls integral on \((\operatorname {E},\mathfrak{T})\) every mapping J of \(\mathcal{M}\) into \(\overline{\mathbf{R}}_+\) satisfying

  the following conditions: \(\alpha)\) \(J(\lambda \cdot f) = \lambda \cdot J(f)\) for \(\lambda \in \overline{\mathbf{R}}_{+}\) and \(f \in \mathcal{M}\) (with the usual convention \(0 \cdot (+\infty) = 0\) ); \(\beta)\) \(J\big(\sum_{n \in \mathbf{N}} f_n\big) = \sum_{n \in \mathbf{N}} J(f_n)\) for every sequence \((f_n)_{n \in \mathbf{N}}\) of elements of \(\mathcal{M}\) .


\begin{enumerate}
\def\labelenumi{\alph{enumi})}
\item
  Let \(J\) be an integral on \((\mathbf{E},\mathfrak{T})\) ; for every subset \(A \in \mathfrak{T}\) , set \(m_{J}(A) = J(\varphi_A)\) ; show that \(m_J\) is an abstract measure and that the mapping \(J \mapsto m_J\) is a bijection of the set of integrals on \((\mathbf{E},\mathfrak{T})\) onto the set of abstract measures on \((\mathbf{E},\mathfrak{T})\) (if \(m\) is an abstract measure, first define \(J(f)\) by linearity for the finite positive \(\mathfrak{T}\) -step functions and use Exer. 3 e)). If \(m\) is an abstract measure, we will denote by \(m^*\) the corresponding integral.
\item
  Let \(m\) be an abstract measure on (E, T). For every positive function \(f\) on E, finite or not, set \(m^{*}(f) = \inf_{g \in \mathcal{M}, g \geqslant f} m^{*}(g)\) . Show that \(m^{*}\) is an encumbrance on E (imitate the proof of Th. 3 of Ch. IV, §1, No.~3). One sets \(m^{*}(A) = m^{*}(\varphi_{A})\) for every subset A of E.
\end{enumerate}

\begin{enumerate}
\def\labelenumi{\arabic{enumi})}
\setcounter{enumi}{4}
\tightlist
\item
  Let \(\mathbf{E}\) be a set, \(\mathfrak{T}\) a tribe of subsets of \(\mathbf{E}\) , \(m\) an abstract measure on \((\mathbf{E},\mathfrak{T})\) , and \(p\geqslant 1\) a finite real number. For every numerical function \(f\) on \(\mathbf{E}\) , one sets \(\mathrm{N}_p(f) = m^* (|f|^p)^{1 / p}\) and one denotes by \(\mathcal{F}^p\) the set of functions \(f\) such that \(\mathrm{N}_p(f)\) is finite (cf.~§1, Exer. 4). One equips the vector space \(\mathcal{F}^p\) with the semi-norm \(\mathrm{N}_p\) , for which it is complete.
\end{enumerate}

\begin{enumerate}
\def\labelenumi{\alph{enumi})}
\item
  Let \(\mathcal{N}\) be the set of functions f such that \(m^{*}(|f|)=0\) , \(\mathcal{V}\) the vector space generated by the finite functions belonging to \(\mathcal{M}\) , and \(\mathcal{L}^{p}\) the smallest closed linear subspace of \(\mathcal{F}^{p}\) containing the characteristic functions of the sets \(A\in \mathfrak{T}\) such that \(m^{*}(A)\) is finite. Show that \(\mathcal{L}^{p}=\mathcal{N}+(\mathcal{V}\cap\mathcal{F}^{p})\) .
\item
  Extend Lebesgue's theorem (Ch. IV, §3, No.~7, Th. 6) to the present case, and examine in particular the case \(p = 1\) (one will define the integral of a function \(f \in \mathcal{L}^1\) ).
\item
  A subset \(\mathbf{A}\) of \(\mathbf{E}\) is said to be \(m\) -integrable if \(\varphi_{\mathbf{A}} \in \mathcal{L}^1\) . Show that there then exists a set \(\mathbf{A}' \in \mathfrak{T}\) such that \(m^{*}(\mathbf{A} \cap \mathbb{C}\mathbf{A}') = m^{*}(\mathbf{A}' \cap \mathbb{C}\mathbf{A}) = 0\) . Converse, in the case that \(m(\mathbf{E})\) is finite.
\end{enumerate}

¶ 6) Let E be a set, T a tribe of subsets of E, and m an abstract measure on (E, T). Suppose that \(m(E)\) is finite and that, for every \(B \in T\) such that \(m(B) > 0\) , there exists a set \(A \in T\) such that \(A \subset B\) and \(0 < m(A) < m(B)\) . For every \(B \in T\) and every number t such that \(0 \leqslant t \leqslant m(B)\) , there exists a set \(A \in T\) such that \(m(A) = t\) and \(A \subset B\) .

¶ 7) Let T be a topological space, \(\Phi\) a clan of subsets of T, and \(m\) a relatively bounded additive mapping (cf.~Exer. 1) of \(\Phi\) into R. Assume that there exists an open covering \(\mathfrak{U}\) of T such that \(\Phi\) consists of the Borel subsets of T contained in the union of a finite number of elements of \(\mathfrak{U}\) . Assume, moreover, that for every \(A \in \Phi\) and every \(\varepsilon > 0\) , there exist a compact subset K of T and an open subset U of T such that \(K \subset A \subset U\) , \(U \in \Phi\) and \(|m|(U - K) < \varepsilon\) . Show that there exists a measure \(\mu\) on T (not necessarily positive) such that \(m(A) = \mu^{\bullet}(A)\) for all \(A \in \Phi\) , and that such a measure \(\mu\) is unique (on introducing the decomposition \(m = m_{+} - m_{-}\) of Exer. 1 c), reduce to the case that \(m\) takes positive values; first treat the case that \(T \in \Phi\) and observe that \(m\) is then inner regular; finally, treat the general case by pasting together measures).

\begin{enumerate}
\def\labelenumi{\arabic{enumi})}
\setcounter{enumi}{7}
\tightlist
\item
  Let T be a topological space and m a bounded, countably additive (cf.~Exer. 1d)) mapping of \(\mathfrak{B}(\mathrm{T})\) into R. Denote by T the set of Borel subsets A of T having the following property: for every \(\varepsilon > 0\) , there exist a closed set F and an open set U in T such that \(F \subset A \subset U\) and \(|m(B)| < \varepsilon\) for every Borel subset B of U - F (in other words, \(|m|(\mathrm{U} - \mathrm{F}) < \varepsilon\) ). Show that T is a tribe of subsets of T (first observe that \(T \in T\) and that T is a clan; it then suffices to prove that if \((\mathrm{A}_{n})_{n \in \mathbb{N}}\) is a sequence of pairwise disjoint elements of T, the set \(A = \bigcup_{n \in N} A_{n}\) belongs to T; to this end, choose closed sets \(F_{n}\) and open sets \(U_{n}\) such that \(F_{n} \subset A_{n} \subset U_{n}\) and \(|m|(\mathrm{U}_{n} - \mathrm{F}_{n}) < \varepsilon / 2^{n}\) , then set \(F = F_{0} \cup \cdots \cup F_{p}\) for p sufficiently large and \(U = \bigcup_{n \in \mathbf{N}} U_{n}\) .
\end{enumerate}


\begin{enumerate}
\def\labelenumi{\arabic{enumi})}
\setcounter{enumi}{8}
\tightlist
\item
  Let \(\mathbf{X}\) be a set; one calls gauge on \(\mathbf{X}\) every countably additive mapping of \(\mathfrak{P}(\mathbf{X})\) into \(\mathbf{R}_+\) ; a gauge \(m\) on \(\mathbf{X}\) is said to be diffuse if \(m(\{x\}) = 0\) for all \(x \in \mathbf{X}\) , and to be atomic if there exists a positive function \(f\) on \(\mathbf{X}\) such that \(m(\mathbf{A}) = \sum_{x \in \mathbf{A}} f(x)\) for every
\end{enumerate}

subset A of X.

\begin{enumerate}
\def\labelenumi{\alph{enumi})}
\item
  Show that every gauge on X may be decomposed in a unique way as the sum of a diffuse gauge and an atomic gauge.
\item
  Let \(X'\) be a subset of \(X\) , and \(m'\) a gauge on \(X'\) ; for every subset \(A\) of \(X\) , one sets \(m(A) = m'(A \cap X')\) . Show that \(m\) is a gauge on \(X\) and that it is diffuse (resp. atomic) if and only if \(m'\) has the same property.
\item
  Let \(m\) be a gauge on \(X\) and let \((X_i)_{i \in I}\) be a family of pairwise disjoint subsets of \(X\) . If every gauge on \(I\) is atomic, then \(m\big(\bigcup_{i \in I} X_i\big) = \sum_{i \in I} m(X_i)\) .
\end{enumerate}

\begin{enumerate}
\def\labelenumi{\arabic{enumi})}
\setcounter{enumi}{9}
\tightlist
\item
  Let \(\mathbf{X}\) be an uncountable infinite set, equipped with a well-ordering relation denoted \(x \leqslant y\) . For every \(x \in \mathbf{X}\) , one denotes by \(\mathrm{I}(x)\) the set of \(y \in \mathbf{X}\) such that \(y < x\) . We make the following hypotheses: \(\alpha\) ) there exists a largest element \(a\) in \(\mathbf{X}\) ; \(\beta\) ) the set of cardinals strictly smaller than \(\operatorname{Card}(\mathbf{X})\) has a largest element \(\mathfrak{c}\) ; \(\gamma\) ) for every \(x\) in \(\mathrm{I}(a)\) , one has \(\operatorname{Card}(\mathrm{I}(x)) \leqslant \mathfrak{c}\) , and the set \(\mathbf{Y}\) of \(x \in \mathbf{X}\) such that \(\operatorname{Card}(\mathrm{I}(x)) = \mathfrak{c}\) is nonempty;
\end{enumerate}

\(\delta)\) every gauge on a set with cardinal \(\leqslant \mathfrak{c}\) is atomic. Show that every gauge on X is atomic. One can argue by contradiction in the following way: let \(m\) be a nonzero diffuse gauge on X. Denote by \(b\) the smallest element of Y, and for every \(x \in I(a)\) let \(f_x\) be an injection of \(I(x)\) into \(I(b)\) . For every pair \((x,y) \in I(a) \times I(b)\) , denote by \(A_{x,y}\) the set of elements \(z \in X\) such that \(x < z < a\) and \(f_z(x) = y\) ; for every \(x \in I(a)\) , denote by \(M_x\) the set of \(y \in I(b)\) such that \(m(A_{x,y}) > 0\) , and for every \(y \in I(b)\) , denote by \(N_y\) the set of \(x \in I(a)\) such that \(m(A_{x,y}) > 0\) . Using Exer. 9 c), show that \(m(X) = \sum_{y \in I(b)} m(A_{x,y})\) ; from this, deduce that the set \(M_x\) is countable and nonempty for every \(x \in I(a)\) ; then show that each of the sets \(N_y\) is countable and from this deduce the contradiction \(\text{Card}(X) = c\) .

\begin{enumerate}
\def\labelenumi{\arabic{enumi})}
\setcounter{enumi}{10}
\tightlist
\item
  A cardinal \(\mathfrak{c}\) is said to be an Ulam cardinal (or to be ulamian) if every gauge on a set with cardinal \(\mathfrak{c}\) is atomic. Show that every countable cardinal is ulamian, that if \(\mathfrak{c}\) is an Ulam cardinal then so is every cardinal \(\mathfrak{c}' \leqslant \mathfrak{c}\) , and that if \((\mathfrak{c}_i)_{i \in I}\) is a family of Ulam cardinals such that \(\operatorname{Card}(I)\) is ulamian, then the cardinal \(\sum_{i \in I} \mathfrak{c}_i\) is ulamian. Finally,
\end{enumerate}

the smallest uncountable cardinal is ulamian (cf.~Exer. 10).

\begin{enumerate}
\def\labelenumi{\arabic{enumi})}
\setcounter{enumi}{11}
\tightlist
\item
  Let X be a set, U a subset of \(\mathfrak{P}(X)\) , and m the characteristic function of U. For m to be a gauge, it is necessary and sufficient that U be an ultrafilter and that the intersection of every countable subset of U belong to U. Under these conditions, U is said to be an Ulam ultrafilter on X. Suppose that this is the case, and that \((\mathrm{H}_{i})_{i\in\mathrm{I}}\) is a family of elements of U such that \(\operatorname{Card}(\mathrm{I})\) is an Ulam cardinal; show that \(\bigcap_{i\in I} H_{i}\) belongs
\end{enumerate}

to \(\mathfrak{U}\) .

\begin{enumerate}
\def\labelenumi{\arabic{enumi})}
\setcounter{enumi}{12}
\tightlist
\item
  In this exercise, we admit the continuum hypothesis, that is, we adjoin to the axioms of set theory the axiom \(\langle\mathrm{Card}(\mathbf{R})\) is the smallest uncountable cardinal \(\rangle\) . \(^{(1)}\)
\end{enumerate}

\begin{enumerate}
\def\labelenumi{\alph{enumi})}
\item
  Every gauge on \(\mathbf{R}\) is atomic (cf.~Exer. 11).
\item
  Let X be a set and m a gauge on X satisfying the following property: for every subset A of X such that \(m(A) > 0\) , there exists a subset B of A such that \(0 < m(B) < m(A)\) ; then m is zero (for every integer \(n \geqslant 1\) , there exists a finite partition \((\mathrm{X}_{n,i})_{i \in \mathrm{I}_n}\) of X such that \(m(\mathrm{X}_{n,i}) \leqslant 1/n\) for all \(i \in I_n\) ; set \(I = \prod_{n \geqslant 1} I_n\) and
\end{enumerate}

\(\mathrm{X}_{\alpha} = \bigcap_{n\geqslant 1}\mathrm{X}_{n,\alpha_n}\) for every \(\alpha = (\alpha_{n})_{n\geqslant 1}\) in I; the family \((\mathrm{X}_i)_{i\in \mathrm{I}}\) is a partition of \(\mathbf{X}\) ,

\(m(X_{i}) = 0\) for every \(i \in I\) , and every gauge on \(I\) is atomic by \(a\) ); conclude by means of Exer. 9 c)).

\begin{enumerate}
\def\labelenumi{\alph{enumi})}
\setcounter{enumi}{2}
\item
  Let X be a set on which there does not exist any nontrivial Ulam ultrafilter; then every gauge on X is atomic (otherwise, by b), there would exist a diffuse gauge m on X and a subset A of X such that \(m(A) > 0\) and such that \(m(B) = 0\) or \(m(A - B) = 0\) for every subset B of A; the set U of subsets B of X such that \(m(A) = m(A \cap B)\) would be a nontrivial Ulam ultrafilter on X).
\item
  Show that \(2^{\mathfrak{c}}\) is an Ulam cardinal for every Ulam cardinal \(\mathfrak{c}\) (by \(c\) ), it suffices to show that if \(C\) is a set with cardinal \(\mathfrak{c}\) , then every Ulam ultrafilter \(\mathfrak{U}\) on \(X = \{0,1\}^{C}\) is trivial; equip \(C\) with a well-ordering structure and define by transfinite induction a family \(s = (s_{\alpha})_{\alpha \in C}\) of elements of \(\{0,1\}\) such that set of \(x = (x_{\alpha})_{\alpha \in C}\) in \(X\) for which \(x_{\alpha}=s_{\alpha}\) for all \(\alpha\leqslant\beta\) belongs to \(\mathfrak{U}\) for every \(\beta\in\mathrm{C}\) (cf.~Exer. 12); then \(s\) belongs to every element of \(\mathfrak{U}\) ). \(^{(2)}\)
\end{enumerate}

\begin{enumerate}
\def\labelenumi{\arabic{enumi})}
\setcounter{enumi}{13}
\item
  Let \((\mathbf{T}_i)_{i\in I}\) be a family of Radon topological spaces, and \(\mathbf{T}\) the sum space of the family. If \(\operatorname{Card}(\mathbf{I})\) is an Ulam cardinal, then the topological space \(\mathbf{T}\) is a Radon space (show that for every positive and bounded countably additive function \(m\) on the Borel tribe of \(\mathbf{T}\) , there exists a countable subset \(J\) of \(I\) such that \(m(\bigcup_{i \in I - J} \mathbf{T}_i) = 0\) .
\item
  Let T be a discrete space whose cardinal is the smallest uncountable cardinal. Show that T is a Radon locally compact space whose topology does not have a countable base (cf.~Exer. 11). Deduce from this that there exist compact Radon spaces that are not metrizable.
\item
  Let \(I\) be a countable set, and \((\mathbf{T}_i)_{i \in I}\) a family of Radon spaces. Assume that every compact subset of any of the spaces \(\mathbf{T}_i\) is metrizable. Show that the product space \(\mathbf{T} = \prod_{i \in I} \mathbf{T}_i\) is Radon. Generalize to the case of inverse limits.
\item
  Let \(\mathbf{T}\) be a topological space.
\end{enumerate}

\begin{enumerate}
\def\labelenumi{\alph{enumi})}
\tightlist
\item
  Let \(m\) be a countably additive mapping of \(\mathfrak{B}(\mathbf{T})\) into \(\mathbf{R}_+\) . Assume that there exists a sequence \((\mathbf{T}_n)_{n\in \mathbf{N}}\) of Borel subsets of \(\mathbf{T}\) such that \(\mathbf{T} = \bigcup_{n\in \mathbf{N}}\mathbf{T}_n\) , \(m(\mathbf{T}_0) = 0\)
\end{enumerate}

and the subspace \(\mathbf{T}_n\) of \(\mathbf{T}\) is Radon for every \(n \geqslant 1\) . Show that there exists a bounded positive measure \(\mu\) on \(\mathbf{T}\) such that \(m(\mathbf{A}) = \mu^{\bullet}(\mathbf{A})\) for every Borel subset \(\mathbf{A}\) of \(\mathbf{T}\) (one can reduce to the case that the \(\mathbf{T}_n\) are pairwise disjoint on applying the Cor. of Prop. 2 of No.~3).

\begin{enumerate}
\def\labelenumi{\alph{enumi})}
\setcounter{enumi}{1}
\item
  Extend the preceding result to the case that the sets \(T_{n}\) for \(n \geqslant 1\) are no longer assumed Borel but merely universally measurable in T (argue as in Prop. 2 of No.~3).
\item
  Every topological space that is the union of a sequence of metrizable compact subspaces is Radon.
\end{enumerate}

\begin{enumerate}
\def\labelenumi{\arabic{enumi})}
\setcounter{enumi}{17}
\tightlist
\item
  Let \(\mathbf{T}_1\) and \(\mathbf{T}_2\) be two topological spaces and \(f\) a bijective continuous mapping of \(\mathbf{T}_1\) onto \(\mathbf{T}_2\) . Assume that, for every bounded countably additive function \(m_1\) on the Borel tribe \(\mathfrak{B}(\mathbf{T}_1)\) , there exists a sequence of compact subsets \(\mathrm{K}_n\) of \(\mathbf{T}_1\) such that \(m_1(\mathbf{T}_1 - \bigcup_{n}\mathbf{K}_n) = 0\) . Let \(m_2\) be a bounded countably additive function on \(\mathfrak{B}(\mathbf{T}_2)\) . In
\end{enumerate}

order that there exist a bounded countably additive function \(m_{1}\) on \(\mathfrak{B}(\mathrm{T}_{1})\) such that \(m_{2}(\mathrm{A}) = m_{1}(f^{-1}(\mathrm{A}))\) for all \(\mathrm{A} \in \mathfrak{B}(\mathrm{T}_{2})\) , it is necessary and sufficient that the following condition be fulfilled: there exists a sequence of compact subsets \(K_{n}\) of \(T_{1}\) such that \(m_{2}(\mathrm{T}_{2} - \bigcup_{n} f(\mathrm{K}_{n})) = 0\) .

\exercisesection{}

\begin{enumerate}
\def\labelenumi{\arabic{enumi})}
\tightlist
\item
  Let \(\mathcal{T} = (\mathrm{K}_i, p_{ij})\) be an inverse system of compact spaces indexed by the set I, let K be a compact space, and let \((p_i)_{i \in I}\) be a coherent and separating family of continuous mappings \(p_i: \mathrm{K} \to \mathrm{K}_i\) .
\end{enumerate}

\begin{enumerate}
\def\labelenumi{\alph{enumi})}
\item
  Let A be the set of continuous functions on K of the form \(f_{i} \circ p_{i}\) , where i runs over I and \(f_{i}\) runs over the set of continuous real functions on \(K_{i}\) . Show that A is a dense linear subspace of \(\mathcal{C}(K)\) .
\item
  Let \((\mu_i)_{i \in I}\) be a sub-inverse system of positive measures on \(\mathcal{T}\) . Assume that the \(p_i\) are surjective. For every \(f \in A\) , let \(I_f\) be the set of \(i \in I\) such that \(f\) is of the form \(f_i \circ p_i\) with \(f_i \in \mathcal{C}(K_i)\) (necessarily unique). Show that there exists one and only one measure \(\pi\) on \(K\) such that \(\pi(f) = \inf_{i \in I_f} \mu_i(f_i)\) for all \(f \in A\) . Show that \(\pi\) is the
\end{enumerate}

greatest of the positive measures \(\mu\) on K such that \(p_i(\mu) \leqslant \mu_i\) for all \(i \in I\) .

¶ 2) Hypotheses and notations are those of Th. 1 of No.~2, for which we propose to indicate a new proof.

\begin{enumerate}
\def\labelenumi{\alph{enumi})}
\item
  Let K be a compact subset of T; for every \(i \in I\) , set \(\mathrm{K}_{i} = p_{i}(\mathrm{T})\) , and denote by \(q_{i}\) the mapping of K onto \(K_{i}\) that coincides on K with \(p_{i}\) ; for \(i \leqslant j\) , denote by \(q_{ij}\) the mapping of \(K_{j}\) into \(K_{i}\) that coincides on \(K_{j}\) with \(p_{ij}\) . Deduce from the preceding exercise the existence of a greatest positive measure \(\pi^{K}\) on K such that \(q_{i}(\pi^{K}) = (\mu_{i})_{\mathbf{K}_{i}}\) for all \(i \in I\) .
\item
  If K and L are compact subsets of T such that \(K \subset L\) , show that \((\pi^{\mathrm{L}})_{\mathrm{K}} \geqslant \pi^{\mathrm{K}}\) ; deduce from this that the set of measures of the form \(i^{\mathrm{K}}(\pi^{\mathrm{K}})\) (where K runs over the set of compact subsets of T, and \(i^{K}\) is the canonical injection of K into T) admits a supremum \(\mu\) in \(\mathcal{M}_{+}(\mathrm{T})\) .
\item
  Show that \(p_i(\mu) = \mu_i\) for every \(i \in I\) (it suffices to observe that \(p_i(\mu) \leqslant \mu_i\) and that the measures \(p_i(\mu)\) and \(\mu_i\) have the same total mass by the hypothesis (P)).
\end{enumerate}

\exercisesection{}

\begin{enumerate}
\def\labelenumi{\arabic{enumi})}
\tightlist
\item
  Let \(\mathbf{T}\) be a topological space and \(\mu\) a positive measure on \(\mathbf{T}\) . Denote by \(\mathfrak{S}_{\mu}\) the set of subsets of \(\mathbf{T}\) whose boundary is \(\mu\) -negligible.
\end{enumerate}

\begin{enumerate}
\def\labelenumi{\alph{enumi})}
\item
  Show that \(\mathfrak{S}_{\mu}\) is a clan.
\item
  If T is completely regular, every point of T has a fundamental system of neighborhoods contained in \(S_{\mu}\) (observe that if f is a continuous function on T, zero outside a \(\mu\) -integrable set, then the set of real numbers r such that \(f(r)\) is not \(\mu\) -negligible is countable).
\end{enumerate}

\begin{enumerate}
\def\labelenumi{\arabic{enumi})}
\setcounter{enumi}{1}
\tightlist
\item
  Let \(\mathbf{T}\) be a completely regular space and \(\mathfrak{F}\) a filter on \(\mathcal{M}_{+}^{b}(\mathrm{T})\) that converges tightly to a bounded measure \(\mu\) . A Borel subset \(\mathbf{A}\) of \(\mathbf{T}\) is said to be a convergence set (for \(\mathfrak{F}\) ) if \(\lim_{\lambda,\mathfrak{F}} \lambda_{\mathbf{A}} = \mu_{\mathbf{A}}\) .
\end{enumerate}

\begin{enumerate}
\def\labelenumi{\alph{enumi})}
\item
  If the disjoint sets \(\mathbf{A}_1\) and \(\mathbf{A}_2\) are convergence sets, then so is \(\mathbf{A}_1 \cup \mathbf{A}_2\) .
\item
  Let A be an open or closed set. For A to be a convergence set, it is necessary and sufficient that \(\lim_{\lambda,\mathfrak{F}} \lambda^{\bullet}(A) = \mu^{\bullet}(A)\) . If A is a convergence set, then so is \(T - A\) .
\item
  Let A be an open or closed convergence set in T, and B a convergence set such that T - B is also a convergence set. Then A ∪ B and A ∩ B are convergence sets, as are their complements.
\item
  The clan generated by the open convergence sets consists of convergence sets.
\item
  Let \(\mathbf{A}\) be a subset of \(\mathbf{T}\) whose boundary is \(\mu\) -negligible. Then \(\mathbf{A}\) is a convergence set.
\item
  Suppose that T is locally compact or Polish. Show that every compact subset K of T is contained in a compact convergence set (in case T is locally compact, use Exercise 1 b). If T is Polish, use the same exercise to construct a sequence of finite coverings \(U_{p}\) of K by open subsets of T, with \(\mu\) -negligible boundary and of diameter \(<2^{-p}\) ; show that \(L=\bigcap_{p}\bigcup_{U\in\mathfrak{U}_{p}}\overline{U}\) is compact and conclude that it is a convergence set by b)).
\end{enumerate}


\begin{enumerate}
\def\labelenumi{\alph{enumi})}
\setcounter{enumi}{6}
\tightlist
\item
  Extend the result of f) to the case of a convergent sequence of measures on a complete metric space.
\end{enumerate}

\begin{enumerate}
\def\labelenumi{\arabic{enumi})}
\setcounter{enumi}{2}
\tightlist
\item
  Let \(\mathbf{T}\) be a Polish space and \((\mu_n)\) a sequence of bounded positive measures on \(\mathbf{T}\) converging tightly to a measure \(\mu\) . Denote by \(\mathfrak{C}\) the set of subsets \(\mathbf{A}\) of \(\mathbf{T}\) having the following property: for every \(\varepsilon > 0\) , there exists a compact subset \(\mathbf{K}\) of \(\mathbf{A}\) such that \(\sup_{n} \mu_n(\mathbf{A} - \mathbf{K}) < \varepsilon\) . Let \(\mathfrak{D}\) be the set of subsets \(\mathbf{A}\) of \(\mathbf{T}\) that, together with their complement, belong to \(\mathfrak{C}\) .
\end{enumerate}

\begin{enumerate}
\def\labelenumi{\alph{enumi})}
\item
  If the sets A and A' belong to C, then so do \(A \cup A'\) and \(A \cap A'\) ; deduce from this that D is a clan of subsets of T.
\item
  Every subset A of T whose boundary is \(\mu\) -negligible belongs to \(\mathfrak{D}\) (apply Prokhorov's theorem to the interior of A, and use Exer. 2 e)).
\item
  Every \(\mathbf{A} \in \mathfrak{D}\) is a convergence set for the sequence \((\mu_n)\) (cf.~Exer. 2); conversely, every open or closed convergence set belongs to \(\mathfrak{D}\) .
\item
  Let \(\mathbf{A} \in \mathfrak{D}\) . Show that there exist a sequence of disjoint compact sets \(\mathbf{K}_p\) and a subset \(\mathbf{N}\) of \(\mathbf{T}\) having the following properties: \(\alpha\) ) each \(\mathbf{K}_p\) is a convergence set for the sequence \((\mu_n); \beta)\) \(\mathbf{A} \subset \mathbf{N} \cup \bigcup_{p} \mathbf{K}_p\) and \(\bigcup_{p} \mathbf{K}_p \subset \mathbf{A} \cup \mathbf{N}; \gamma)\) for every \(\varepsilon > 0\) , there exists an open neighborhood \(\mathbf{U}\) of \(\mathbf{N}\) such that \(\sup_{n}\mu_n(\mathbf{U}) < \varepsilon\) (apply Exer. 2 f) and Prokhorov's theorem to a suitable subspace of \(\mathbf{T}\) that is the intersection of a sequence of open sets).
\end{enumerate}

\begin{enumerate}
\def\labelenumi{\arabic{enumi})}
\setcounter{enumi}{3}
\item
  Let \(\mathbf{T}\) be a completely regular space, \(t\) a point of \(\mathbf{T}\) , and \(\mathfrak{U}\) a set of Borel subsets of \(\mathbf{T}\) that generates the filter of neighborhoods of \(t\) . Let \(\mathbf{I}\) be a set equipped with a filter \(\mathfrak{F}\) , and let \((\mu_i)_{i \in \mathbb{I}}\) be a family of bounded positive measures of total mass 1 on \(\mathbf{T}\) . In order that \(\lim_{i, \mathfrak{F}} \mu_i = \varepsilon_t\) , it is necessary and sufficient that \(\lim_{i, \mathfrak{F}} \mu_i(\mathrm{U}) = 1\) for every \(\mathrm{U} \in \mathfrak{U}\) .
\item
  Let E be a real Hilbert space admitting an orthonormal basis \((x_{n})_{n\in \mathbf{N}}\) . Let T be the space E equipped with the weakened topology, and let \((a_{n})_{n\in \mathbf{N}}\) be a sequence of real numbers such that \(0 < a_{n} < 1\) for all \(n\) and \(\lim_{n\to \infty}n^2\log a_n = 0\) . For every \(n\in \mathbf{N}\) , set \(\mu_{n} = \sum_{p\in \mathbf{N}}a_{n}^{p}(1 - a_{n})\cdot \varepsilon_{n\cdot x_{p}}\) . Show that \(\mu_{n}\) is a bounded positive measure on T of total mass 1 for every \(n \in \mathbf{N}\) , that \(\lim_{n\to\infty}\mu_{n}=\varepsilon_{0}\) (tight convergence), and that \(\lim_{n\to\infty}\mu_{n}(\mathrm{K})=0\) for every compact subset \(\mathrm{K}\) of T (apply the criterion of the preceding exercise to the set \(\mathfrak{U}\) of neighborhoods of 0 of the form \(\{x \mid (a \mid x) \leqslant 1\}\), where a runs over E). In particular, the set of elements of the sequence \((\mu_{n})_{n \in \mathbf{N}}\) is a relatively compact subset of \(\mathcal{M}_{+}^{b}(\mathbf{T})\) that does not satisfy Prokhorov's condition.
\end{enumerate}


\begin{enumerate}
\def\labelenumi{\arabic{enumi})}
\setcounter{enumi}{5}
\tightlist
\item
  Let \(\mathbf{T}\) be a topological space and \(\mathbf{H}\) a set of positive measures on \(\mathbf{T}\) ; assume that the following condition is fulfilled:
\end{enumerate}

\begin{enumerate}
\def\labelenumi{(\Alph{enumi})}
\setcounter{enumi}{21}
\tightlist
\item
  Every point of \(\mathbf{T}\) admits an open neighborhood \(\mathbf{W}\) such that \(\sup_{\mu \in \mathbf{H}} \mu^{\bullet}(\mathbf{W})\) is finite.
\end{enumerate}

Finally, let \(\mathfrak{U}\) be an ultrafilter on H.

\begin{enumerate}
\def\labelenumi{\alph{enumi})}
\item
  Let \(\mathbf{K}\) be a compact subset of \(\mathbf{T}\) . Show that the induced measures \(\mu_{\mathbf{K}} (\mu \in \mathbf{H})\) converge vaguely with respect to \(\mathfrak{U}\) to a measure \(\pi^{\mathbf{K}}\) on \(\mathbf{K}\) .
\item
  Let \(\mathbf{K}\) and \(\mathbf{L}\) be two compact subsets of \(\mathbf{T}\) such that \(\mathbf{K} \subset \mathbf{L}\) ; show that \((\pi^{\mathbf{L}})_{\mathbf{K}} \geqslant \pi^{\mathbf{K}}\) .
\item
  For every compact subset K of T, let \(i^{\mathbf{K}}\) be the canonical injection of K into T. Show that the family of measures \(i^{\mathbf{K}}(\pi^{\mathbf{K}})\) admits a supremum \(\pi\) in \(\mathcal{M}_{+}(\mathrm{T})\) .
\item
  Let \(f\) be a lower semi-continuous positive function on \(\mathbf{T}\) , and \(g\) an upper semi-continuous positive function on \(\mathbf{T}\) with compact support. Show that \(\pi^{\bullet}(f) \leqslant \lim_{\mu, \mathfrak{U}} \mu^{\bullet}(f)\) and \(\pi^{\bullet}(g) \geqslant \lim_{\mu, \mathfrak{U}} \mu^{\bullet}(g)\) .
\end{enumerate}

¶ 7) Let H be a set of bounded positive measures on a topological space T. Assume that \(\sup_{\mu\in H}\mu^{\bullet}(T)\) is finite, and that for every \(\varepsilon>0\) , there exists a compact subset K of T

such that \(\sup_{\mu \in \mathbf{H}}\mu^{\bullet}(\mathrm{T} - \mathrm{K}) < \varepsilon\) . Show that the condition (V) of the preceding exercise is

verified. Let \(\mathfrak{U}\) be an ultrafilter on H. The measure \(\pi\) is defined as in the preceding exercise.

\begin{enumerate}
\def\labelenumi{\alph{enumi})}
\item
  Show that \(\pi^{\bullet}(g) \geqslant \lim_{\mu, \mathfrak{U}} \mu^{\bullet}(g)\) for every upper semi-continuous bounded positive function \(g\) on T. Deduce from this that \(\pi^{\bullet}(f) = \lim_{\mu, \mathfrak{U}} \mu^{\bullet}(f)\) for every bounded function \(f\) on T whose set of points of discontinuity is \(\pi\) -negligible.
\item
  When T is completely regular, deduce from a) that H is relatively compact for the tight topology (this furnishes a new proof of Th. 1 of No.~5 for the case of positive measures).
\end{enumerate}

¶ 8) Let T be a complete separable metric space, and d its metric. For every closed subset F of T and every real number a > 0, denote by \(F^{a}\) the set of \(x \in T\) such that \(d(x, F) < a\) . Given two bounded positive measures \(\lambda\) and \(\mu\) on T, denote by \(D(\lambda, \mu)\) the infimum of the set of real numbers a > 0 satisfying the inequalities \(\lambda^{\bullet}(F) \leqslant \mu^{\bullet}(F^{a}) + a\) , \(\mu^{\bullet}(F) \leqslant \lambda^{\bullet}(F^{a}) + a\) for every closed subset F of T.

\begin{enumerate}
\def\labelenumi{\alph{enumi})}
\item
  Show that D is a metric on \(\mathcal{M}_+^b\) and that \(\mathrm{D}(\varepsilon_x,\varepsilon_y) = d(x,y)\) for two points \(x\) and \(y\) of T such that \(d(x,y) < 1\) .
\item
  Let \(f\) be a Borel mapping of \(T\) into \(T\) (cf.~TG, IX, §6, No.~3 or GT, IX, §6, Exer. 16) and \(\lambda\) a bounded positive measure on \(T\) . Show that there exists a bounded positive measure \(\mu\) on T such that \(\mu^{\bullet}(\mathrm{A}) = \lambda^{\bullet}\big(f^{-1}(\mathrm{A})\big)\) for every Borel subset A of T. Let \(a > 0\) , and let H be the set of \(x \in \mathrm{T}\) such that \(d(x, f(x)) \geqslant a\) ; show that H is Borel in T and that \(\mathrm{D}(\lambda, \mu) \leqslant \sup (a, \lambda(\mathrm{H}))\) (for every closed subset F of T, one has \(\mathrm{F} \cap (\mathrm{T} - \mathrm{H}) \subset f(\mathrm{F}^a)\) and \(\lambda^{\bullet}(\mathrm{F}) \leqslant \lambda^{\bullet}(\mathrm{F} \cap (\mathrm{T} - \mathrm{H})) + \lambda^{\bullet}(\mathrm{H})\) ).
\item
  Let \(\mu\) and \(\mu_n\) (for \(n \geqslant 1\) ) be bounded positive measures on \(\mathbf{T}\) such that \(\lim_{n \to \infty} \mathrm{D}(\mu_n, \mu) = 0\) . Show that the sequence \((\mu_n)\) converges tightly to \(\mu\) (show that \(\lim_{n \to \infty} \mu_n^\bullet(\mathrm{T}) = \mu^\bullet(\mathrm{T})\) and \(\mu^\bullet(\mathrm{F}) \geqslant \limsup_{n \to \infty} \mu_n^\bullet(\mathrm{F})\) for every closed subset \(\mathbf{F}\) of \(\mathbf{T}\) ; deduce from this that \(\mu^\bullet(f) \leqslant \liminf_{n \to \infty} \mu_n^\bullet(f)\) for every lower semi-continuous function \(f \geqslant 0\) by the method of Lemma 3 of §2, No.~6).
\item
  Conversely, show that for every sequence of bounded positive measures \(\mu_n\) tending tightly to \(\mu\) in \(\mathcal{M}_+^b\) , one has \(\lim_{n \to \infty} D(\mu_n, \mu) = 0\) . One can proceed as follows: let \(\varepsilon > 0\) , and let \(K\) be a compact subset of \(T\) such that \(\mu^\bullet(T - K) < \varepsilon\) and \(\sup_{n} \mu_n^\bullet(T - K) < \varepsilon\) ;
\end{enumerate}

construct a finite family \((\mathrm{B}_{i})_{1\leqslant i\leqslant p}\) of sets with \(\mu\) -negligible boundary, of diameter \(\leqslant\varepsilon/2\) , pairwise disjoint, whose union contains K (cf.~Exer. 1). Let f be a mapping of T into T, constant on each of the sets \(B_{1},\ldots,B_{p}\) and \(\mathbb{C}(B_{1}\cup\cdots\cup B_{p})\) and such that \(f(\mathrm{B}_i) \subset \mathrm{B}_i\) for \(1 \leqslant i \leqslant p\) . Define the measures \(\pi_n\) and \(\pi\) by \(\pi_n^\bullet(\mathrm{A}) = \mu_n^\bullet(\overline{f}(\mathrm{A}))\) and \(\pi^\bullet(\mathrm{A}) = \mu^\bullet(f^{-1}(\mathrm{A}))\) for every Borel subset A of T; deduce from b) the relations \(\mathrm{D}(\pi_n, \mu_n) \leqslant \varepsilon\) , \(\mathrm{D}(\pi, \mu) \leqslant \varepsilon\) and show that \(\lim_{n \to \infty} \mathrm{D}(\pi_n, \pi) = 0\) .

\begin{enumerate}
\def\labelenumi{\alph{enumi})}
\setcounter{enumi}{4}
\tightlist
\item
  The metric D on \(M_{+}^{b}\) is compatible with the tight topology (observe that \(M_{+}^{b}\) is metrizable for the tight topology).
\end{enumerate}

¶ 9) Notations and hypotheses are those of the preceding exercise. Let \((\mu_n)\) be a Cauchy sequence in \(\mathcal{M}_+^b\) for the metric D; show that the sequence \((\mu_n)\) is convergent (this gives a new proof of the fact that the space \(\mathcal{M}_+^b\) is Polish for the tight topology). One can proceed thus:

\begin{enumerate}
\def\labelenumi{\alph{enumi})}
\item
  Let \(\varepsilon > 0\) and \(a > 0\) be two real numbers; there exists a finite subset \(F\) of \(T\) such that \(\sup_{n} \mu_n^\bullet(T - F^a) \leqslant \varepsilon\) (choose an integer \(N \geqslant 1\) and a compact subset \(K\) of \(T\) such that \(\sup_{n \geqslant N} D(\mu_n, \mu_N) \leqslant \varepsilon/2\) and \(\mu_N^\bullet(T - K) \leqslant \varepsilon/2\) ; deduce from this that \(\sup_{n \geqslant N} |\mu_n^\bullet(T) - \mu^\bullet(T)| \leqslant \varepsilon\) and \(\sup_{n \geqslant N} \mu_n^\bullet(T - K^{a/2}) \leqslant \varepsilon\) ; finally, choose a finite subset \(F\) of \(T\) such that \(K^{a/2} \subset F^a\) and \(\sup_{n \geqslant N} \mu_n^\bullet(T - F^a) \leqslant \varepsilon\) ).
\item
  Let \(\varepsilon > 0\) ; there exists a compact subset \(K\) of \(T\) such that \(\sup_{n} \mu_n^\bullet(T - K) \leqslant \varepsilon\) (for every integer \(p \geqslant 1\) choose a finite subset \(F_p\) of \(T\) such that \(\mu_n^\bullet(T - (F_p)^{2^{-p}}) \leqslant \varepsilon / 2^p\) , and set \(K = \bigcap_{p\geqslant 1}(\overline{(F_p)^{2^{-p}}})\) .
\item
  Deduce from \(b\) ) that the Cauchy sequence \((\mu_n)\) has a convergent subsequence (for the metric D).
\end{enumerate}

¶ 10) let T be a completely regular space; denote by S the set of real functions \(f \geqslant 1\) on T, such that the set of points t of T for which \(f(t) \leqslant c\) is compact for every real number c.~For every \(f \in S\) , denote by \(\mathcal{M}_f\) the set of bounded measures \(\mu\) on T such that \(|\mu|^{\bullet}(f) \leqslant 1\) .

\begin{enumerate}
\def\labelenumi{\alph{enumi})}
\item
  For a subset H of \(\mathcal{M}^{b}(T)\) to satisfy Prokhorov's condition, it is necessary and sufficient that there exist an \(f \in S\) such that \(H \subset M_{f}\) (for necessity, take f to be of the form \(c(1 + \sum_{n \geqslant 1} n \cdot f_{n})\) where \(f_{n}\) is the characteristic function of a set \(U_{n}\) such that \(\mathrm{T} - \mathrm{U}_n\) is compact and \(\sup_{\mu \in \mathbf{H}} |\mu|^{\bullet} (\mathrm{U}_n) \leqslant 2^{-n}\) .
\item
  Suppose T is locally compact. For a subset H of \(\mathcal{M}^{b}(T)\) to be relatively compact (for the tight topology), it is necessary and sufficient that there exist a continuous function \(f \in S\) such that \(H \subset M_{f}\) .
\item
  Let C be a closed convex cone in \(\mathcal{M}_{+}^{b}(\mathrm{T})\) . Show that \(\mathrm{C} \cap \mathcal{M}_{f}\) is a cap of C (TVS, II, §7, No.~2, Def. 3) for every \(f \in \mathcal{S}\) , and deduce from this that C is the union of its caps. Denote by E the union of the extremal generators of C. Suppose T is Souslin. Show that for every \(\pi \in \mathrm{C}\) , there exist a Borel subset B of \(\mathcal{M}_{+}^{b}(\mathrm{T})\) contained in E and a positive measure P on B of total mass 1, such that \(\pi = \int_{\mathrm{B}} \mu d\mathrm{P}(\mu)\) (apply Choquet's integral representation theorem).
\end{enumerate}

\begin{enumerate}
\def\labelenumi{\arabic{enumi})}
\setcounter{enumi}{10}
\item
  Let T be a completely regular space. Assume given, for every compact space K and every continuous mapping f of T into K, a positive measure \(\mu_{f,K}\) on K; assume that \(g(\mu_{f,K}) = \mu_{g \circ f,L}\) for any compact spaces K and L and any continuous mappings \(f : T \to K\) and \(g : K \to L\) ; assume moreover that the measure \(\mu_{f,K}\) is concentrated on \(f(T)\) for any f and K. Show that there exists one and only one bounded positive measure \(\pi\) on T such that \(\mu_{f,K} = f(\pi)\) for any f and K (make use of the universal property of the Stone--Čech compactification of T).
\item
  Let T be a completely regular space, a subspace of a locally compact space L. For every measure \(\mu\) (positive or not) on T, there exists a locally compact subspace \(T'\) of L containing T, and a measure \(\mu'\) on \(T'\) concentrated on T that induces \(\mu\) on T.
\end{enumerate}

¶ 13) *Let G be a locally compact abelian group. One denotes by \(\widehat{G}\) the dual group of G and by \(\alpha\) a Haar measure on \(\widehat{G}\) . For every bounded measure \(\mu\) on G, one denotes by \(\mathcal{F}\mu\) the Fourier transform of \(\mu\) , which is a function on the group \(\widehat{G}\) (cf.~Theor. spec., Ch. II, §1, No.~2).

\begin{enumerate}
\def\labelenumi{\alph{enumi})}
\item
  The Fourier transformation \(\mathcal{F}\) is a continuous mapping of \(\mathcal{M}_{+}^{b}(\mathrm{G})\) equipped with the tight topology into \(\mathcal{C}^b (\widehat{\mathrm{G}})\) equipped with the topology of compact convergence. (Make use of the Cor. of Prop. 13 of No.~6.)
\item
  Let \(\mathfrak{F}\) be a filter on \(\mathcal{M}_+^b (\mathrm{G})\) and \(\Phi\) a bounded continuous function on \(\widehat{\mathrm{G}}\) ; assume that \(\lim_{\lambda ,\mathfrak{F}}(\mathcal{F}\lambda)\cdot \alpha = \Phi \cdot \alpha\) (vague convergence) and \(\lim_{\lambda ,\mathfrak{F}}(\mathcal{F}\lambda)(0) = \Phi (0)\) . Show that the filter \(\mathfrak{F}\) converges tightly to a measure \(\mu\) such that \(\mathcal{F}\mu = \Phi\) (show by application of the Stone-Weierstrass theorem that the set of Fourier transforms of continuous functions on G with compact support is a dense linear subspace of \(\mathcal{C}^0 (\overline{\mathrm{G}})\) for uniform convergence; next observe that \(\lim_{\lambda ,\mathfrak{F}}\int (\mathcal{F}\lambda)\cdot u d\alpha = \int \Phi u d\alpha\) for every \(u\in \mathrm{E}\) and that the filter \(\mathfrak{F}\) contains a vaguely compact set of bounded measures; deduce from this that the filter \(\mathfrak{F}\) has at most one cluster point, then that it converges tightly).
\item
  Let \(\mathfrak{F}\) be a filter on \(\mathcal{M}_{+}^{b}(\mathrm{G})\) having a countable base. In order that \(\mathfrak{F}\) converge tightly to a bounded positive measure \(\mu\) , it is necessary and sufficient that there exist a continuous function \(\Phi\) on \(\widehat{\mathbf{G}}\) such that \(\mathcal{F}\lambda\) converges pointwise to \(\Phi\) with respect to \(\mathfrak{F}\) , in which case \(\mathcal{F}\mu = \Phi\) ( \(P\) . Lévy's theorem).*
\end{enumerate}

\exercisesection{}

¶ 1) For every compact interval K of R, let \(\mathcal{C}(K)\) be the real vector space of continuous functions on K, equipped with the topology of uniform convergence; denote by \(\mathcal{D}(K)\) the vector space of infinitely differentiable real functions on R that are zero outside K. Equip \(\mathcal{D}(K)\) with the coarsest topology for which the mappings \(f \mapsto D^{p}f|K\) are continuous for every positive integer \(p\) (D is the differentiation operator). Set \(\mathcal{D} = \bigcup_{K} \mathcal{D}(K)\) , a space that one equips with the locally convex topology that is the direct limit

of the topologies of the subspaces \(\mathcal{D}(\mathbf{K})\) .

\begin{enumerate}
\def\labelenumi{\alph{enumi})}
\tightlist
\item
  For every integer \(p \geqslant 0\) and every function \(f \in \mathcal{D}(\mathbf{K})\) , set
\end{enumerate}

\[
\mathrm{Q} _ {p} (f) = \int_ {\mathrm{K}} \left(\mathrm{D} ^ {p} f (x)\right) ^ {2} d x.
\]

Show that \(\mathbf{Q}_p\) is a positive quadratic form on \(\mathcal{D}(\mathbf{K})\) , that the sequence of norms \(\mathbf{Q}_p^{1/2}\) defines the topology of \(\mathcal{D}(\mathbf{K})\) , and that \(\operatorname{Tr}(\mathbf{Q}_{p+1}/\mathbf{Q}_p)\) is finite for every \(p \geqslant 0\) . (For every real \(t\) , let \(\mathrm{I}_t\) be the characteristic function of the interval \([t, +\infty[\) ; using the formula \(\mathrm{D}^p f(t) = \int_{\mathbf{K}} \mathrm{D}^{p+1} f \cdot \mathrm{I}_t dx\) and Bessel's inequality, prove that \(\sum_{i=1}^{n} \mathbf{Q}_p(f_i) \leqslant l^2\) for every finite sequence \(f_1, \ldots, f_n\) of functions belonging to \(\mathcal{D}(\mathbf{K})\) , orthonormal for \(\mathbf{Q}_{p+1}\) ; the number \(l\) is the length of the interval K.) Deduce from this that \(\mathcal{D}(\mathbf{K})\) is a nuclear space.

\begin{enumerate}
\def\labelenumi{\alph{enumi})}
\setcounter{enumi}{1}
\tightlist
\item
  Prove that the space \(\mathcal{D}\) is nuclear. (First establish the existence of a nonzero function in \(\mathcal{D}\) , and deduce from this the existence of a function \(h \geqslant 0\) in \(\mathcal{D}\) such that \(\sum_{n \in \mathbf{Z}} h(x - n) = 1\) for every \(x \in \mathbf{R}\) . Let \(K\) be a compact interval of \(\mathbf{R}\) such that \(h\) is
\end{enumerate}

zero outside of K; for every integer \(n\) , set \(h_n(x) = h(x - n)\) and \(\mathrm{K}_n = \mathrm{K} + n\) . Let V be a convex neighborhood of 0 in \(\mathcal{D}\) ; there exist positive integers \(p_n\) such that every function \(f \in \mathcal{D}(\mathrm{K}_n)\) with \(\mathrm{Q}_{p_n}(f) \leqslant 1\) belongs to V. Define the continuous quadratic forms Q and R on \(\mathcal{D}\) by \(\mathrm{Q}(f) = \sum_{n \in \mathbf{Z}} 2^{2n} \mathrm{Q}_{p_n}(f \cdot h_n)\) and \(\mathrm{R}(f) = \sum_{n \in \mathbf{Z}} 2^{3n} \mathrm{Q}_{p_n + 1}(f \cdot h_n)\) ; show that V contains the set of \(f \in \mathcal{D}\) such that \(\mathrm{Q}(f) \leqslant 1\) and that \(\operatorname{Tr}(\mathrm{R}/\mathrm{Q})\) is finite.)

\(^{*}\) c) Generalize the foregoing to infinitely differentiable functions on \(R^{n}\) with compact support. \(^{*}\)

\specsection{Exercises for the Annex}

\begin{enumerate}
\def\labelenumi{\arabic{enumi})}
\tightlist
\item
  Let E be a finite-dimensional vector space over a commutative field of characteristic \(\neq 2\) . Denote by H a nondegenerate quadratic form on E, and by S the symmetric bilinear form on \(E \times E\) such that \(\mathrm{H}(x) = \mathrm{S}(x, x)\) for all x in E. Let Q be a quadratic form on E.
\end{enumerate}

\begin{enumerate}
\def\labelenumi{\alph{enumi})}
\item
  There exists an endomorphism \(u\) of \(\mathbf{E}\) characterized by the relations \(\mathbf{Q}(x) = \mathrm{S}(u(x), x)\) and \(\mathrm{S}(u(x), y) = \mathrm{S}(x, u(y))\) for all \(x\) and \(y\) in \(\mathbf{E}\) . One sets \(\operatorname{Tr}(\mathbf{Q}/\mathbf{H}) = \operatorname{Tr}(u)\) .
\item
  Generalize Remark 3 of No.~1.
\item
  If \((e_i)_{1\leqslant i\leqslant m}\) is a basis of E orthonormal for H, then \(\operatorname{Tr}(\mathbf{Q}/\mathbf{H}) = \sum_{i=1}^{m}\mathrm{Q}(e_i)\) .
\end{enumerate}

\begin{enumerate}
\def\labelenumi{\arabic{enumi})}
\setcounter{enumi}{1}
\tightlist
\item
  Let E be a real Hilbert space, Q a continuous positive quadratic form on E, and H the quadratic form \(x \mapsto \|x\|^2\) on E.
\end{enumerate}

\begin{enumerate}
\def\labelenumi{\alph{enumi})}
\item
  Show that \(\operatorname{Tr}(\mathbf{Q}/\mathbf{H}) = \sum_{i \in I} \mathbf{Q}(e_i)\) for every orthonormal basis \((e_i)_{i \in I}\) of E. (Let \((a_1, \ldots, a_p)\) be a finite orthonormal family in E; for every \(\varepsilon > 0\) , there exist a finite subset J of I and elements \(a_1', \ldots, a_p'\) that are linear combinations of the \(e_i\) for \(i \in J\) and are such that \(\|a_j - a_j' \| < \varepsilon\) for \(1 \leqslant j \leqslant p\) ; then \(\sum_{j=1}^{p} \mathbf{Q}(a_j') \leqslant \sum_{i \in J} \mathbf{Q}(e_j) \leqslant \sum_{i \in I} \mathbf{Q}(e_i)\) . Deduce from this that \(\sum_{j=1}^{p} \mathbf{Q}(a_j) \leqslant \sum_{i \in I} \mathbf{Q}(e_i)\) .)
\item
  Deduce from \(a\) ) a new proof of Prop. 3.
\item
  Let \(\mathbf{E}_0\) be a dense linear subspace of \(\mathbf{E}\) , \(\mathbf{Q}_0\) (resp. \(\mathbf{H}_0\) ) the restriction of \(\mathbf{Q}\) (resp. \(\mathbf{H}\) ) to \(\mathbf{E}_0\) . Show that \(\operatorname{Tr}(\mathbf{Q}/\mathbf{H}) = \operatorname{Tr}(\mathbf{Q}_0/\mathbf{H}_0)\) . (Let \((e_1, \ldots, e_n)\) be a linearly independent sequence in \(\mathbf{E}\) , generating a subspace \(\mathbf{F}\) . Denote by \(\mathbf{Q}_{\mathbf{F}}\) (resp. \(\mathbf{H}_{\mathbf{F}}\) ) the restriction of \(\mathbf{Q}\) (resp. \(\mathbf{H}\) ) to \(\mathbf{F}\) . Using Remark 3 of No.~1, show that \(\operatorname{Tr}(\mathbf{Q}_{\mathbf{F}}/\mathbf{H}_{\mathbf{F}})\) is a continuous function of \((e_1, \ldots, e_n)\) .)
\end{enumerate}

\begin{enumerate}
\def\labelenumi{\arabic{enumi})}
\setcounter{enumi}{2}
\item
  Let E and F be two real vector spaces and u a linear mapping of E into F; if Q and H are positive quadratic forms on F such that \(\mathrm{H}(x)=0\) implies \(\mathrm{Q}(x)=0\) for \(x\in F\) , then \(\mathrm{Tr}\left(\mathrm{Q}\circ\mathrm{u}/\mathrm{H}\circ\mathrm{u}\right)\leqslant\mathrm{Tr}\left(\mathrm{Q}/\mathrm{H}\right)\) .
\item
  Let I be a set and \(l^2(\mathbf{I})\) the vector space of families \(\mathbf{x} = (x_i)_{i \in \mathbf{I}}\) of real numbers such that \(\sum_{i \in \mathbf{I}} x_i^2\) is finite. Let \((\lambda_i)_{i \in \mathbf{I}}\) be a summable family of positive numbers. For every \(\mathbf{x}\) in \(l^2(\mathbf{I})\) , set \(Q(\mathbf{x}) = \sum_{i \in \mathbf{I}} \lambda_i x_i^2\) and \(H(\mathbf{x}) = \sum_{i \in \mathbf{I}} x_i^2\) . Show that \(\operatorname{Tr}(Q / H) = \sum_{i \in \mathbf{I}} \lambda_i\) . (Make use of Exercise 2 a).
\item
  Let E be a real vector space, \((\lambda_i)_{i \in I}\) a summable family of positive numbers, and \((y_i)_{i \in I}\) a family of linear forms on E. For every \(x \in E\) , set \(\mathrm{H}(x) = \sum_{i \in I} \langle x, y_i \rangle^2\) and \(\mathrm{Q}(x) = \sum_{i \in I} \lambda_i \langle x, y_i \rangle^2\) ; assume that \(\mathrm{H}(x)\) is finite for every \(x \in E\) . Show that \(\mathrm{Q}(x)\) is finite for every \(x \in E\) , that Q and H are positive quadratic forms on E, and that \(\operatorname{Tr}(\mathrm{Q}/\mathrm{H}) \leqslant \sum_{i \in I} \lambda_i\) . (Set \(u(x) = (\langle x, y_i \rangle)_{i \in I}\) and apply Exer. 3 to the linear mapping \(u\) of E into \(l_2(I)\) .)
\item
  Let E be a real vector space, Q, H and \(H_{0}\) positive quadratic forms on E. Assume that \(H \leqslant a \cdot H_{0}\) , where a is a positive real number. Prove that \(\operatorname{Tr}(Q/H_{0}) \leqslant a \cdot \operatorname{Tr}(Q/H)\) . (Reduce to the case that E is finite-dimensional by Remark 1 of No.~1, then conclude by Prop. 1, using the existence of a basis of E that is orthogonal for both H and \(H_{0}\) .)
\item
  Let E be a real Hilbert space. Show that the Sazonov topology on E is defined by the set of semi-norms \(Q^{1/2}\) , where Q runs over the set of nuclear positive quadratic forms on E. (Use Exer. 6.)
\item
  Let E and F be two real Hilbert spaces. There exists on \(E \otimes F\) a Hausdorff pre-Hilbert space structure such that \((x_{1} \otimes y_{1}|x_{2} \otimes y_{2}) = (x_{1}|x_{2}) \cdot (y_{1}|y_{2})\) for \(x_{1}, x_{2}\) in E and \(y_{1}, y_{2}\) in F. One denotes by \(E \otimes_{2} F\) the Hilbert space completion of \(E \otimes F\) .
\end{enumerate}

\begin{enumerate}
\def\labelenumi{\alph{enumi})}
\item
  If \(\mathbf{E}'\) (resp. \(\mathbf{F}'\) ) is a closed linear subspace of \(\mathbf{E}\) (resp. \(\mathbf{F}\) ), show that \(\mathbf{E}' \otimes_2 \mathbf{F}'\) may be identified with the closed linear subspace of \(\mathbf{E} \otimes_2 \mathbf{F}\) generated by the elements \(x \otimes y\) such that \(x \in \mathbf{E}'\) and \(y \in \mathbf{F}'\) .
\item
  Suppose that E (resp. F) is the hilbertian sum of a family \((\mathrm{E}_{\alpha})_{\alpha\in\mathbf{A}}\) (resp. \((\mathrm{F}_{\beta})_{\beta\in\mathbf{B}}\) ) of closed linear subspaces (TVS, V, §2, No.~2, Def. 2). Show that \(E\otimes_{2}F\) is the hilbertian sum of the family \((\mathrm{E}_{\alpha}\otimes_{2}\mathrm{F}_{\beta})_{(\alpha,\beta)\in\mathbf{A}\times\mathbf{B}}\) of closed linear subspaces.
\item
  If \((e_i)_{i\in \mathbf{I}}\) (resp. \((f_j)_{j\in \mathbf{J}})\) is an orthonormal basis of E (resp. F), then the family \((e_i\otimes f_j)_{(i,j)\in \mathbf{I}\times \mathbf{J}}\) is an orthonormal basis of \(\mathbf{E}\otimes_2\mathbf{F}\) .
\item
  Let \(\mathbf{G}\) be a real Hilbert space. Define canonical isomorphisms of \(\mathbf{E} \otimes_{2} \mathbf{F}\) onto \(\mathbf{F} \otimes_{2} \mathbf{E}\) and of \((\mathbf{E} \otimes_{2} \mathbf{F}) \otimes_{2} \mathbf{G}\) onto \(\mathbf{E} \otimes_{2} (\mathbf{F} \otimes_{2} \mathbf{G})\) .
\item
  Let \(E_{1}\) and \(F_{1}\) be two real Hilbert spaces, and \(u:E\to E_{1}, v:F\to F_{1}\) continuous linear mappings. Show that \(u\otimes v\) may be extended by continuity to a continuous linear mapping \(u\otimes_{2}v\) of \(E\otimes_{2}F\) into \(E_{1}\otimes_{2}F_{1}\) , and that \(\|u\otimes_{2}v\|=\|u\|\cdot\|v\|\) .
\end{enumerate}

\begin{enumerate}
\def\labelenumi{\arabic{enumi})}
\setcounter{enumi}{8}
\tightlist
\item
  Let E and F be two real Hilbert spaces.
\end{enumerate}

\begin{enumerate}
\def\labelenumi{\alph{enumi})}
\item
  Show that there exists a continuous linear mapping \(\varphi\) of \(\mathbf{E} \otimes_2 \mathbf{F}\) into \(\mathcal{L}(\mathbf{E}; \mathbf{F})\) characterized by \(\varphi(x \otimes y)(x') = (x|x') \cdot y\) for \(x, x'\) in \(\mathbf{E}\) and \(y\) in \(\mathbf{F}\) . Show that \(\varphi\) has norm 1, and that \(\varphi(\mathbf{E} \otimes \mathbf{F})\) is the set of continuous linear mappings of finite rank of \(\mathbf{E}\) into \(\mathbf{F}\) .
\item
  Show that \(\varphi\) is a bijection of \(E \otimes_2 F\) onto the set HS (E, F) of Hilbert--Schmidt linear mappings of E into F (make use of Exer. 8 c) and Prop. 3). Equip HS (E, F) with the Hilbert space structure deduced from that of \(E \otimes_2 F\) by means of the bijection \(\varphi\) ; the corresponding norm is denoted \(\|u\|_2\) . Let \(u \in \mathrm{HS}(E, F)\) ; one defines positive quadratic forms \(Q_u\) and H on E by \(Q_u(x) = \|u(x)\|^2\) and \(H(x) = \|x\|^2\) . Show that \(\|u\|_2^2 = \operatorname{Tr}(Q_u / H)\) .
\item
  Let \(\mathbf{E}_1\) and \(\mathbf{F}_1\) be two real Hilbert spaces, and \(u: \mathbf{E}_1 \to \mathbf{E}\) , \(v: \mathbf{F} \to \mathbf{F}_1\) continuous linear mappings. Let \(\varphi_1\) be the isomorphism of \(\mathbf{E}_1 \otimes_2 \mathbf{F}_1\) onto \(\mathrm{HS}(\mathbf{E}_1, \mathbf{F}_1)\) defined in a manner analogous to \(\varphi\) . Show that \(v \circ \varphi(t) \circ u = \varphi_1((u^* \otimes v)(t))\) for all \(t \in \mathbf{E} \otimes_2 \mathbf{F}\) . From this, deduce that for every Hilbert-Schmidt mapping \(w\) of \(\mathbf{E}\) into \(\mathbf{F}\) , the linear mapping \(v \circ w \circ u\) of \(\mathbf{E}_1\) into \(\mathbf{F}_1\) is Hilbert-Schmidt, and that \(\|v \circ w \circ u\|_2 \leqslant \|v\| \cdot \|w\|_2 \cdot \|u\|\) .
\end{enumerate}

\specialchapter{Historical Note}

(N.B. --- The Roman numerals refer to the bibliography at the end of this note.)

While the study of the connections between topology and measure theory goes back to the beginnings of the modern theory of functions of real variables, it is only very recently that integration in Hausdorff topological spaces has been thoroughly worked out in general fashion. Before setting out the history of the work that preceded the present synthesis, we review several stages in the evolution of the ideas regarding the relations between topology and measure.

For Lebesgue, it is only a question of integrating functions of one or several real variables. In 1913, Radon defined general measures on \(\mathbf{R}^n\) and the corresponding integrals; this theory is exposed in detail in the book (I) of Ch. de la Vallée Poussin and rests, in constant fashion, on the topological properties of Euclidean spaces. A little later, in 1915, Fréchet defined in (II, \(a\) ) `abstract' measures on a set equipped with a tribe, and the integrals with respect to these measures; he noted that one can establish in this way the principal results of the Lebesgue theory without the use of topological methods. He justified his undertaking with the following words, taken from the introduction of (II, \(b\) ), first part): « Que par exemple dans l'espace à une infinité de coordonnées où diverses applications de l'Analyse avaient conduit à diverses définitions non équivalentes d'une suite convergente, on remplace une de ces définitions par une autre, rien ne sera changé dans les propriétés des familles et fonctions additives d'ensembles dans ces espaces».(*) Fréchet's investigations were completed by Carathéodory, to whom is due an important theorem on the extension of a set function to a measure. The beginning of the book of Saks (III) gives a condensed exposition of this point of view.

The discovery of Haar measure on locally compact groups (cf.~the Historical Note for Chs. VII and VIII) and the numerous applications that it immediately received, then the works of Weil and Gelfand on Harmonic Analysis, led around 1940 to a profound modification of this point of view: in this kind of question, it is most convenient to regard a measure as a linear form on a space of continuous functions. This method obliges one to restrict oneself to compact or locally compact spaces, but this is not an inconvenience for nearly all of the applications: better yet, the introduction of Harmonic Analysis on p-adic groups and adèle groups by J. Tate and A. Weil has permitted a spectacular renewal of Analytic Number Theory.

It is from an entirely different direction that the need arises for enlarging this point of view by considering measures on non locally compact topological spaces: gradually, Probability Theory leads to the study of such spaces and furnishes numerous nontrivial examples. Perhaps the reason for the tardy influence of these developments on the theory of measure is to be found in the relative isolation of Probability Theory, which remained until recently on the fringes of the traditional mathematical disciplines.

\specsection{Measures on spaces of sequences}

One of the most highly developed branches of classical Probability Theory is that of limit theorems (law of large numbers, tendency towards the Gauss--Laplace law, \ldots); this concerns a deepening of the concept of statistical regularity manifested by phenomena that bring into play very large populations. The correct mathematical formulation of these problems requires the introduction of measures on sequence spaces; these spaces, which constitute the most obvious generalization of finite-dimensional spaces, were the subject of choice of investigations into `General Analysis' undertaken around 1920 by Fréchet, Lévy, Lusin, \ldots{} Nor is it fortuitous that Khintchine and Kolmogoroff, the creators of the new methods of Probability Theory, were both disciples of Lusin, and that Lévy very quickly oriented himself towards probabilistic problems: these constituted the touchstone of the new methods.

The first implicit intervention of a measure on a sequence space appeared in the work devoted by E. Borel in 1909 to countable probabilities (IV). A highly original idea of Borel consists in the application of the probabilistic results he had just obtained to the proof of properties possessed by the decimal expansion of almost every real number between 0 and 1. This application rests on the following fundamental remark: let us define every real number between 0 and 1 by the sequence of figures (or `digits') in its expansion in a given base \(q\) ( \(q \geqslant 2\) ); if one successively draws at random the various figures of a number \(x\) , independently of each other and with equal probability \(1 / q\) for \(0,1,\ldots ,q - 1\) , the probability that \(x\) lies in an interval of {[}0,1{[} is equal to the length of that interval.

In 1923, Steinhaus (V) established these results rigorously and described the precise mathematical model for the infinite sequence of random drawings considered by Borel: for simplicity let us take q=2 and denote by I the set with two elements \(\{0,1\}\) ; one equips I with the measure \(\mu\) defined by \(\mu(0)=\mu(1)=\frac{1}{2}\) ; the elements of the product space \(I^{N}\) are the sequences \(\varepsilon=(\varepsilon(n))_{n\in\mathbf{N}}\) of numbers equal to 0 or 1, and the mapping \(\varphi:\varepsilon\mapsto\sum_{n\geqslant0}\varepsilon(n)\cdot2^{-n-1}\) is, up to a countable set, a bijection of \(I^{N}\) onto the interval {[}0,1{]}; moreover, \(\varphi^{-1}\) transforms the Lebesgue measure on {[}0,1{]} into the measure P on \(I^{N}\) that is the product of the measures \(\mu\) on each of the factors. Actually, Steinhaus did not have at his disposal a construction for product measures; he used the existence of the quasi-bijection \(\varphi\) to construct the measure P on \(I^{N}\) starting from Lebesgue measure on {[}0,1{]}, and then gave an axiomatic characterization of P. The isomorphism so obtained made it possible to translate the language of probability into that of measure and to apply the known theorems on the Lebesgue integral.

In the same work, Steinhaus considered the random series \(\sum_{n\geqslant 0}\sigma_n\cdot a_n\) where the signs \(\sigma_{n} = \pm 1\) are chosen at random independently of each other and with equal probability \(\frac{1}{2}\) ; between 1928 and 1935, he studied numerous other random series. From their side, Paley, Wiener and Zygmund considered random Fourier series(1) of the form \(\sum_{n = -\infty}^{\infty}a_n\exp (2\pi i(nt + \Phi_n))\) ; the `amplitudes' \(a_{n}\) are fixed, and the `phases' \(\Phi_{n}\) are independent random variables uniformly distributed on {[}0,1{]}. While the analytic difficulties vary enormously from one of these problems to another, the translation in terms of measure theory is the same in all cases and represents an extension of the case treated by Borel and Steinhaus; it is a matter of constructing a measure on \(\mathbf{R}^{\mathbf{N}}\) that is the product of a family of measures all identical to a same positive measure \(\mu\) on \(\mathbf{R}\) of mass 1; for example, the preceding random Fourier series correspond to the case that \(\mu\) is the Lebesgue measure on {[}0,1{]}.

To construct such product measures, two methods can be used. The first is a direct method, precisely exposed for the first time by Daniell (VI, a)) in 1918; it was rediscovered in 1934 by Jessen (VII), who made a detailed study of the case that \(\mu\) is Lebesgue measure on {[}0, 1{]}. The second method is to seek artifices analogous to that of Steinhaus to reduce to Lebesgue measure on {[}0, 1{]}; this way of proceeding had the advantage of convenience so long as one did not have available a complete exposition of general measure theory, for it permitted using Lebesgue's theorems without the need for new proofs. \(^{(2)}\)

\specsection{The theory of Brownian motion}

This theory occupies a special position in contemporary scientific development, by the constant and fertile exchange between physical problems and `pure' mathematics to which it bears witness. The study of Brownian motion, discovered in 1829 by the botanist Brown, had been conducted intensively in the 19th century by numerous physicists, \(^{(3)}\) but the first satisfactory mathematical model was only invented by Einstein in 1905. In the simple case of a particle moving along a straight line, Einstein's fundamental hypotheses may be formulated as follows: if \(x(t)\) is the abscissa of the particle at the instant t, and if \(t_{0} < t_{1} < \cdots < t_{n-1} < t_{n}\) , the successive displacements \(x(t_{i}) - x(t_{i-1})\) (for \(1 \leqslant i \leqslant n\) ) are independent Gaussian random variables. This is not the place to evoke in detail the important experimental work of J. Perrin that motivated Einstein's theory; for our purposes, we need only retain a remark of Perrin, according to which the observation of trajectories of Brownian motion irresistibly suggested to him the ``mathematicians' functions without derivative''. This remark was to be the initial spark for Wiener.

Quite another current of ideas has its origins in the kinetic theory of gases, developed between 1870 and 1900 by Boltzmann and Gibbs. Consider a gas formed by N molecules of mass m at (absolute) temperature T, and denote by \(v_{1},\ldots,v_{N}\) the velocities of the N molecules of the gas; the kinetic energy of the system is equal to

\[
\frac {m}{2} (\mathbf {v} _ {1} ^ {2} + \dots + \mathbf {v} _ {\mathrm{N}} ^ {2}) = 3 \mathrm{N} k \mathrm{T},\tag{1}
\]

where k is Boltzmann's constant. According to the ideas of Gibbs, the multitude of shocks between molecules does not allow the precise determination of the velocities of the molecules, and it is convenient to introduce a probability law P on the sphere S of the space of dimension 3N defined by the equation (1). The `micro-canonical' hypothesis consists in assuming that P is the measure on the sphere S with mass 1 invariant under rotation. Moreover, Maxwell's law of velocities states that the law of probability of a component of the velocity of a molecule is a Gaussian measure with variance \(2k\mathrm{T} / m\) (§6, No.~5, Remark 3). Borel seems to have been the first to observe in 1914 that Maxwell's law is a consequence of Gibbs' hypotheses and the properties of the sphere when the number of molecules is very large. He considered a sphere S in a Euclidean space of large dimension and the measure P of mass 1 on S invariant under rotation; using the classical approximation methods based on Stirling's formula, he showed that the projection of P on a coordinate axis is approximately Gaussian. These results were sharpened a little later by Gâteaux and Lévy (IX, \(a\) ). Given an integer \(m \geqslant 1\) and a number \(r > 0\) , denote by \(\mathrm{S}_{m,r}\) the set of sequences of the form \((x_{1}, \ldots, x_{m}, 0, 0, \ldots)\) with \(x_{1}^{2} + \cdots + x_{m}^{2} = r^{2}\) ; also, denote by \(\sigma_{m,r}\) the measure with mass 1 on \(\mathrm{S}_{m,r}\) invariant under rotation. Stated in modern language, the result of Gâteaux and Lévy is as follows: the sequence of measures \(\sigma_{m,1}\) tends tightly to a unit mass at the origin \((0, 0, \ldots)\) , and the sequence of measures \(\sigma_{m,\sqrt{m}}\) tends tightly to a measure \(\Gamma\) of the form

\[
d \Gamma (x _ {1}, x _ {2}, \dots) = \prod_ {n = 1} ^ {\infty} d \gamma (x _ {n})
\]

(γ is the Gaussian measure on R with variance 1).

The preceding measure \(\Gamma\) plays the role of a Gaussian measure in infinite dimensions. It seems indeed that Lévy had confusedly hoped to define in an intrinsic manner a Gaussian measure on every infinite-dimensional Hilbert space. In fact, as shown by Lévy and Wiener, the measure \(\Gamma\) is invariant in a certain sense \(^{(4)}\) under the automorphisms of \(l^{2}\) ; unfortunately, the set \(l^{2}\) of square-summable sequences \((x_{1}, x_{2}, \ldots, x_{n}, \ldots)\) has measure zero for \(\Gamma\) . It is now known that one must be content to have a Gaussian promeasure on an infinite-dimensional Hilbert space. \(^{(5)}\)

We owe to Wiener the essential progress: if one does not have a reasonable Gaussian measure on an infinite-dimensional Hilbert space, one can construct by the operation of primitive a measure \(w\) on a space of continuous functions starting from a Gaussian promeasure (cf.~§6, No.~7, Th. 1 for the details). We shall explain succinctly Wiener's original construction of \(w\) (X); it is directly influenced by the relation \(\Gamma = \lim_{m \to \infty} \sigma_{m,\sqrt{m}}\) of Gâteaux and Lévy. For every integer \(m \geqslant 1\) , denote by \(H_m\) the set of functions on \(T = ]0,1]\) that are constant on each of the intervals \(\left[ \frac{k-1}{m}, \frac{k}{m} \right]\) (for \(k = 1,2,\ldots,m\) ), and by \(\pi_m\) the measure of mass 1, invariant under rotation, on the Euclidean sphere of radius 1 in \(R^m\) . Let \(f_m\) be the isomorphism of \(H_m\) onto \(R^m\) that associates, to each function taking the value \(a_k\) on the interval \(\left[ \frac{k-1}{m}, \frac{k}{m} \right]\) , the vector \((a_1, a_2 - a_1, \ldots, a_m - a_{m-1})\) (whence the term `differential space' dear to Wiener); denote by \(w_m\) the measure on \(H_m\) that is the image of \(\pi_m\) under \(f_m^{-1}\) . Wiener defined the desired measure \(w\) as the limit of the measures \(w_m\) . To be precise, let us denote by \(H\) the set of regulated functions on \(T\) , with the topology of uniform convergence (we have \(H_m \subset H\) for every integer \(m \geqslant 1\) ); for every bounded uniformly continuous function \(F\) on \(H\) , the limit \(A\{F\} = \lim_{m \to \infty} \int_{H_m} F(x) dw_m(x)\) exists; next, Wiener obtained certain upper bounds by a subtle analysis of the fluctuations of the game of heads-or-tails, and taking up again the compactness arguments highlighted by Daniell, he showed that one is under the conditions for applying Daniell's extension theorem. One concludes the existence of a measure \(w\) carried by \(C(T)\) and such that \(A\{F\} = \int_{C(T)} F(x) dw(x)\) . Wiener was then able to show that the measure \(w\) corresponds to Einstein's hypotheses, \(^{(6)}\) and his estimates allowed him to give a precise meaning to Perrin's remark on functions without derivatives: the set of functions satisfying a Lipschitz condition of order \(\frac{1}{2}\) is negligible for \(w\) (however, for every \(a\) with \(0 < a < \frac{1}{2}\) , almost every function satisfies a Lipschitz condition of order \(a\) ).

Today, numerous constructions of the Wiener measure are known. Thus, Paley and Wiener use random Fourier series (XI, Ch. IX): for every real

\[
\begin{array}{l} \int_ {\mathcal {C} (\mathrm{T})} f (x (t _ {1}), \dots , x (t _ {n})) d w (x) = \\ (2 \pi) ^ {- n / 2} \prod_ {i = 1} ^ {n} (t _ {i} - t _ {i - 1}) ^ {- 1 / 2} \int \dots \int f (x _ {1}, \dots , x _ {n}) \exp \left(- \frac {1}{2} \sum_ {i = 1} ^ {n} \frac {(x _ {i} - x _ {i - 1}) ^ {2}}{t _ {i} - t _ {i - 1}}\right) d x _ {1} \dots d x _ {n}, \end{array}
\]

sequence \(\mathbf{a} = (a_n)_{n\geqslant 1}\) and every integer \(m\geqslant 0\) , let us define the function \(f_{m,\mathbf{a}}\) on ]0,1] by

\[
f _ {m, \mathbf {a}} (t) = a _ {1} t + 2 \sum_ {k = 2} ^ {2 ^ {m + 1}} \frac {1}{\pi k} a _ {k - 1} \sin \pi k t;
\]

one can show that for \(\Gamma\) -almost every sequence \(\mathbf{a}\) , the sequence of functions \(f_{m,\mathbf{a}}\) tends to a continuous function \(f_{\mathbf{a}}\) , and that \(w\) is the image of \(\Gamma\) under the mapping (defined almost everywhere) \(\mathbf{a} \mapsto f_{\mathbf{a}}\) . Later on, Lévy gave in (IX, \(b\) ), \(c\) ) a construction very near to that which we have exposed in §6, No.~7. Finally, Kac, Donsker and Erdős showed around 1950 how to replace the spherical measures \(\pi_m\) on \(\mathbf{R}^m\) in Wiener's original construction by more general measures. Their results establish a solid link between Wiener measure and the limit theorems of Probability Theory; they were to be completed and systematized by Prokhorov (XIII) in a work to which we shall return later on.

This is not the place to analyze the numerous and important probabilistic works brought about by Wiener's discovery; nowadays, Brownian motion appears only as one of the most important examples of a Markoff process. We shall only mention Kac's application of Wiener measure to the solution of certain parabolic partial differential equations; this is a question of adapting ideas of Feynmann in quantum field theory---yet another example of the reciprocal influence of mathematics and problems of physics.

\specsection{Inverse limits of measures}

This is a theory that has developed mainly in response to the needs of Probability Theory. Problems concerning a finite sequence of random variables \(X_{1},\ldots,X_{n}\) are in principle solved when one knows the law \(P_{X}\) of the sequence: this is a positive measure on \(R^{n}\) of mass 1 such that the probability of obtaining simultaneously the inequalities \(a_{1}\leqslant X_{1}\leqslant b_{1},\ldots,a_{n}\leqslant X_{n}\leqslant b_{n}\) is equal to \(\mathrm{P}_{\mathrm{X}}(\mathrm{C})\) , where C is the closed box \([a_{1},b_{1}]\times\cdots\times[a_{n},b_{n}]\) in \(R^{n}\) . In practice, the measure \(P_{X}\) either has discrete support or admits a density with respect to Lebesgue measure. When dealing with an infinite sequence \((\mathrm{X}_{n})_{n\geqslant1}\) of random variables, one generally knows the law \(P_{n}\) of the partial sequence \((\mathrm{X}_{1},\ldots,\mathrm{X}_{n})\) for every integer \(n\geqslant1\) ; these data satisfy a compatibility condition that expresses that the sequence \((\mathrm{P}_{n})_{n\geqslant1}\) is an inverse system (or `projective system') of measures. Until about 1920, the probabilities of events relating to the infinite sequence were defined in more or less implicit fashion by `natural' passages to the limit from the probabilities of the finite case; one thus assumes that the probability that a game terminates is the limit, as n tends to infinity, of the probability that it terminates in at most n steps. Naturally, such a theory is not very coherent, and nothing excludes the presence of `paradoxes', the same probability receiving two distinct estimates according as one evaluates by means of one or the other of two procedures, each as `natural' as the other.

Steinhaus (V) appears to have been the first to have felt the need for considering (for the game of heads-or-tails) not only the inverse system \((\mathbf{P}_{n})_{n\geqslant1}\) but its limit as well. A little earlier, in 1919, Daniell (VI, b)) had proved in general the existence of such inverse limits, \(^{(7)}\) but this result seems to have remained unknown in Europe. It was rediscovered in 1933 by Kolmogoroff in the work (XII) where the author formulates the axiomatic conception of Probability Theory. The proofs of Daniell and Kolmogoroff make use of a compactness argument, which is substantially the same as the one we have employed in Th. 2 of §4, No.~3 and which rests on Dini's theorem.

The Daniell--Kolmogoroff theorem left nothing to be desired for the case of random sequences \((\mathbf{X}_{n})_{n\geqslant1}\) , but the study of random functions undertaken from 1935 on by Kolmogoroff, Feller and Doob harbors difficulties of quite another order. Consider for example an interval T of R, representing the set of instants of observation of a `stochastic process'; the set of possible trajectories is the product space \(R^{T}\) , regarded as the inverse limit of the partial products \(R^{H}\) , where H runs over the set of finite subsets of T; one generally assumes given an inverse system of measures \((\mu_{\mathrm{H}})\) (cf.~§4, No.~2). Kolmogoroff's theorem does indeed yield a measure on \(R^{T}\) , but it is defined on a tribe notably smaller than the Borel tribe. \(^{(8)}\) A variant of Kolmogoroff's construction, which yields a measure on a topological space, is due to Kakutani (Proc. Imp. Acad. Tokyo 19 (1943), 184--188), and has been rediscovered several times since: one regards \(\mu_{H}\) as a measure on \(\overline{R}^{H}\) carried by \(R^{H}\) ; \(^{(9)}\) the compact space \(E = \overline{R}^{T}\) is the inverse limit of the finite products \(\overline{R}^{H}\) and one can define a measure \(\mu\) on E as the inverse limit of the \(\mu_{H}\) (cf.~Ch. III, §4, No.~5). However, this procedure has a serious inconvenience; the elements of \(\overline{R}^{T}\) possess no regularity property that permits advancing the probabilistic study of the process---or even simply eliminating the parasitic values \(\pm\infty\) introduced by the compactification \(\overline{R}\) of R. This can be remedied by inducing the measure \(\mu\) of \(\overline{R}^{T}\) on a particular subspace (for example \(\mathcal{C}(T)\) in the case of Brownian motion); the fundamental difficulty arises from the fact that a function space, even one of the usual type, is not necessarily \(\mu\) -measurable in \(\overline{R}^{T}\) , and even the choice of function space may be questionable. \(^{(10)}\)

A decisive step was taken in 1956 by Prokhorov in a work (XIII) that has had a determining influence on the theory of stochastic processes. By putting into a suitable axiomatic form the methods used by Wiener in the article analyzed above, he established a general existence theorem for inverse limits of measures on function spaces that is a special case of Th. 1 of §4, No.~2 corresponding to Polish spaces.

A more restricted class of inverse systems was introduced by Bochner (XIV) in 1947; these are inverse systems formed by finite-dimensional real vector spaces and surjective linear mappings. The inverse limit of such a system may be identified in a natural way with the algebraic dual \(\mathbf{E}^*\) of a real vector space \(\mathbf{E}\) , equipped with the weak topology \(\sigma (\mathbf{E}^{*},\mathbf{E})\) ; a corresponding inverse system of measures has a limit that is a measure \(\mu\) defined on a tribe notably smaller than the Borel tribe of \(\mathbf{E}^*\) . Bochner completely characterized such `promasures' by their Fourier transform, which is a function on \(\mathbf{E}\) . But this result is scarcely usable in the absence of a topology on \(\mathbf{E}\) , in which case one has to examine the possibility of regarding \(\mu\) as a measure on the topological dual \(\mathbf{E}'\) of \(\mathbf{E}\) . In an independent way, R. Fortet and E. Mourier, while seeking to generalize to random variables with values in a Banach space certain classical results of Probability Theory (law of large numbers, central limit theorem) also brought into evidence the fundamental role played by the Fourier transformation in such questions. But substantial progress was not made until 1956 when Gelfand (XV, \(b\) ) suggested that the natural setting for the Fourier transformation is not that of Banach spaces or Hilbert spaces, but that of nuclear Fréchet spaces. He conjectured that every continuous function of positive type on such a space is the Fourier transform of a measure on its dual, a result established soon afterwards by Minlos (XVI). Its importance stems above all from the fact that it is applicable to spaces of distributions, and that the quasi-totality of function spaces are Borel subsets of the space of distributions (which thus constitutes a much better receptacle than \(\mathbf{R}^{\mathrm{T}}\) ). \(^{(11)}\) The theory of random distributions is an area in full expansion, and we shall content ourselves by referring the reader to the book of Gelfand and Vilenkin (XVII).

The results on inverse limits just mentioned make use of the existence of topologies on the base spaces. One may ask whether there exists an analogous theory in the case of `abstract' measures. Von Neumann proved in 1935 the existence of product measures in all cases, but the discovery of a counter-example by Jessen and Andersen (XVIII) dashed the hope that every inverse system of measures admits a limit. Two palliatives have been discovered: in 1949, C. Ionescu-Tulcea established the existence of countable inverse limits, by means of the existence of suitable disintegrations, \(^{(12)}\) a highly interesting result for the study of Markoff processes; moreover, it had been observed that the topology of the spaces only intervenes through the intermediary of the set of compact subsets. It was therefore natural to try to axiomatize this situation within the abstract theory, by means of the concept of compact class of subsets of a set. This work was done in 1953 by Marczewski (who established an abstract inverse limit theorem by such means) and Ryll-Nardzewski (who treated the disintegration of measures). \(^{(13)}\)

\specsection{Measures on general topological spaces and tight convergence}

The study of the connections between topology and measure theory has above all been conceived of as the study of regularity properties of measures, and in particular that of `outer' regularity and of `inner' regularity; \(^{(14)}\) inner regularity is equivalent to outer regularity on a locally compact space countable at infinity. The construction that Lebesgue gives to the measure of subsets of the line highlights these two kinds of regularity, and the outer regularity property of measures on a Polish space seems to have had public notoriety around 1935. But it was not until 1940, in an article whose diffusion was retarded by the war, that A. D. Alexandroff (XIX) gave prominence to the role of inner regularity and showed that it is possessed by the measures on a Polish space; this result was rediscovered later by Prokhorov (XIII) and is often erroneously attributed to this author. It was not perceived until very recently that this property extends to Souslin spaces; because of this, the importance of these spaces has increased greatly, even more so as it was realized that their theory could be worked out without any metrizability hypothesis, and that the quasi-totality of the function spaces were Souslin (most often, even Lusin). \(^{(15)}\) These are the reasons that moved us to place the accent on inner regular measures in this chapter.

The definition of a mode of convergence (vague or tight) for measures is most conveniently done by putting the space of measures into duality with a space of continuous functions. Generalizing an old result of F. Riesz, A. A. Markoff established in 1938 a one-to-one correspondence between the positive functionals on \(\mathcal{C}(\mathrm{X})\) and the regular measures on a compact space X. In the work (XIX) already cited, A. D. Alexandroff extends these results to the case of a completely regular space: he introduces a hierarchy in the set of positive linear forms on the space \(\mathcal{C}^b (\mathrm{X})\) of bounded continuous functions on a completely regular space \(\mathrm{X},^{(16)}\) he defines tight convergence of bounded measures and proves among others the following two theorems:

\begin{enumerate}
\def\labelenumi{\alph{enumi})}
\item
  if X is Polish, the set of linear forms on \(\mathcal{C}^{b}(X)\) corresponding to the measures is closed for the weak convergence of sequences;
\item
  if a sequence of bounded measures has a tight limit, `no mass escapes at infinity' (this is a weak form of the converse of Prokhorov's theorem on tight convergence).
\end{enumerate}

From this abundance of concepts and theorems, Prokhorov was able to extract the results important for the theory of stochastic processes, and to present them in a simple and striking form. In his great work of 1956 already cited (XIII), a large part is devoted to bounded positive measures on a Polish space; generalizing a construction of Lévy, he defines a metric on the set of positive measures of mass 1 that makes it a Polish space, and then establishes an important compactness criterion for tight convergence (cf.~§5, No.~5, Th. 1). Independently of Prokhorov, Le Cam (XX) obtained a number of compactness results for the tight convergence of measures; he makes no metrizability hypothesis on the spaces he considers, and his results reduce to earlier theorems of Dieudonné in the locally compact case.

\specsection{Bibliography}

\begin{enumerate}
\def\labelenumi{(\Roman{enumi})}
\item
  Ch. DE LA VALLÉE POUSSIN, Intégrales de Lebesgue, Fonctions d'ensembles, Classes de Baire, Paris (Gauthier--Villars), 1st. edn., 1916; 2nd. edn., 1936.
\item
  M. FRÉCHET: a) Sur l'intégrale d'une fonctionelle étendue à un ensemble abstrait, Bull. Soc. Math. France, 43 (1915), pp.~248--265; b) Les familles et fonctions additives d'ensembles abstraits, Fund. Math., 4 (1923), pp.~329--365; ibid. 5 (1924), pp.~206-251.
\item
  S. SAKS, Theory of the Integral, 2nd. edn., New York (Stechert), 1937.
\item
  E. BOREL, Les probabilités dénombrables et leurs applications arithmétiques, Rend. Circ. Math. Palermo, 27 (1909), pp.~247-271.
\item
  H. STEINHAUS, Les probabilités dénombrables et leur rapport à la théorie de la mesure, Fund. Math., 4 (1923), pp.~286--310.
\item
  P. J. DANIELL: a) Integrals in an infinite number of dimensions, Ann. of Math., 20 (1918--19), pp.~281--288; b) Functions of limited variation in an infinite number of dimensions, ibid. 21 (1919--20), pp.~30--38.
\item
  B. JESSEN, The theory of integration in a space of an infinite number of dimensions, Acta Math., 63 (1934), pp.~249-323.
\item
  A. EINSTEIN, Investigations on the Theory of the Brownian Movement, New York (Dover), 1956.
\item
  P. LÉVY: a) Leçons d'Analyse Fonctionnelle, Paris (Gauthier-Villars), 1922 (the second edition appeared in 1951 with the same publisher, under the title Problèmes concrets d'Analyse Fonctionnelle); b) Processus stochastiques et mouvement brownien, Paris (Gauthier-Villars), 1948; c) Le mouvement brownien, Mémor. Sci. Math., 126 (1954).
\item
  N. WIENER, Differential space, J. Math. Phys. MIT, 2 (1923), pp.~131--174 (= Selecta, pp.~55--98, Cambridge (MIT Press), 1964).
\item
  R. E. A. C. PALEY and N. WIENER, Fourier transforms in the complex domain, Amer. Math. Soc. Colloq. Publ. No.~19, New York, 1934.
\item
  A. N. KOLMOGOROFF, Grundbegriffe der Wahrscheinlichkeit-srechnung, Berlin (Springer), 1933.
\item
  Ju. V. PROKHOROV, Convergence of random processes and limit theorems in probability theory, Theor. Probability Appl., 1 (1956), pp.~156--214.
\item
  S. BOCHNER, Harmonic Analysis and the Theory of Probability, Berkeley (University of California Press), 1960.
\item
  I. M. GELFAND: a) Generalized stochastic processes (in Russian), Dokl. Akad. Nauk SSSR, 100 (1955), pp.~853--856; b) Some problems of functional analysis (in Russian), Uspehi Mat. Nauk, 11 (1956), pp.~3--12 (= Amer. Math. Soc. Transl. (2), 16 (1960), pp.~315--324).
\item
  R. A. MINLOS, Generalized random processes and their extension to a measure (in Russian), Trudy Moskov. Math. Obšč., 8 (1959), pp.~497--518 (= Selected translations in mathematical statistics and probability, III (19), pp.~291--313).
\item
  I. M. GELFAND and N. Ya. VILENKIN, Generalized Functions, Vol. IV, New York (Academic Press), 1964 (English translation).
\item
  E. SPARRE-ANDERSSEN and B. JESSEN, On the introduction of measures in infinite product sets, Dansk. Vid. Selbskab. Mat. Fys. Medd., 25 (1948), No.~4, pp.~1-7.
\item
  A. D. ALEXANDROFF, Additive set functions in abstract spaces. I (Ch. 1), Mat. Sbornik, 8 (1940), pp.~307--348; II (Chs. 2 and 3), ibid., 9 (1941), pp.~563--628; III (Chs. 4 to 6), ibid., 13 (1943), pp.~169--238.
\item
  L. LE CAM, Convergence in distribution of stochastic processes, Univ. Cal. Publ. Statistics, No.~11 (1957), pp.~207--236.
\end{enumerate}

\specialchapter{Index of notations}

Reference numbers indicate, in order, the chapter, section and subsection.

Chapter VII:

\(\gamma_{X}(s)\) , \(\gamma(s)\) : VII, 1, 1.

\(\gamma(s)f\) , \(\gamma(s)\mu\) (f a function, \(\mu\) a measure): VII, 1, 1.

\(d\mu(s^{-1}x)\) : VII, 1, 1.

\(\delta_{X}(s)\) , \(\delta(s)\) , \(\delta(s)f\) , \(\delta(s)\mu\) , \(d\mu(xs)\) : VII, 1, 1.

\(\check{f}\) , \(\check{\mu}\) , \(d\mu(x^{-1})\) (f a function, \(\mu\) a measure): VII, 1, 1.

\(\Delta_{G}\) , \(\Delta\) : VII, 1, 3.

mod \(_{G}\varphi\) , mod \(\varphi\) ( \(\varphi\) an automorphism): VII, 1, 4.

\(Z_{p}\) (p a prime number): VII, 1, 6.

\(K^{+}\) (K a field): VII, 1, 10.

mod \(_{K}a\) , mod a (a an element of a locally compact field K): VII, 1, 10.

\(\mathcal{H}^{\chi}(X)\) , \(\mathcal{H}_{+}^{\chi}(X)\) , \(\mathcal{H}^{1}(X)\) , \(f^{\chi}\) , \(f^{1}\) (X a locally compact space in which a locally compact group H operates, \(\chi\) a continuous representation of H in \(R_{+}^{*}\) : VII, 2, 1.

\(f^{b}\) : VII, 2, 2.

\(\lambda^{\sharp}\) , \(\frac{\mu}{\beta}\) , \(\mu/\beta\) : VII, 2, 2.

\(m^{\sharp}\) (m a vectorial measure): VII, 2, 2.

\(T_{J}\) , \(T_{1}(n,K)\) , \(T(n,K)\) , \(T(n,K)^{*}\) : VII, 3, 3.

Chapter VIII :

\(^{n}_{i=1}\mu_{i}\) , \(*_{\varphi}(\mu_{i})_{1\leqslant i\leqslant n}\) , \(\mu_{1}*\mu_{2}*\cdots*\mu_{n}\) : VIII, 1, 1.

\(\gamma_{\chi}\) : VIII, 2, 3 and VIII, 2, 4.

\(\gamma_{\chi,p}\) : VIII, 2, 5.

\(U(\mu)\) (U a representation of a locally compact group G, \(\mu\) a measure on G): VIII, 2, 6.

\(\mathcal{M}^{\rho}(G)\) (G a locally compact group): VIII, 3, 1.

\(\mu*\beta f\) , \(\mu*f\) ( \(\mu\) a measure, f a function): VIII, 4, 1.

\(\mathcal{L}(G)\) (G a locally compact group): VIII, 4, 5.

\(\mathcal{U}_{s}^{\infty}(G)\) (G a locally compact group): VIII, 4, Exer. 21.

Chapter IX:

\(\mathcal{F}_{+}(\mathrm{T}), \mathcal{F}_{+}, f_{\mathrm{A}}, f^{0}\): preliminary conventions.

\(\pi(p), p_{\mathrm{A}}\) or \(p|_{\mathrm{A}}: \mathrm{IX}, 1, 1\).

\(\mathcal{P}(\mathrm{T}; \mathbf{C}), \mathcal{P}(\mathrm{T}; \mathbf{R}), \mathcal{P}(\mathrm{T}), \mathcal{P}_{+}(\mathrm{T}): \mathrm{IX}, 1, 2\).

\(w^{\bullet}(f), \int^{\bullet} f dw, \int^{\bullet} f(t) dw(t): \mathrm{IX}, 1, 2\).

\(w^{\bullet}, w_{\mathrm{K}}^{\bullet}: \mathrm{IX}, 1, 2\).

\(w^{+}, w^{-}, |w|: \mathrm{IX}, 1, 2\).

\(\mu(f), \mu(\mathrm{A}): \mathrm{IX}, 1, 5\).

\(\operatorname{Supp}(\mu): \mathrm{IX}, 1, 6\).

\(\sum_{i \in \mathrm{I}} \mu_i: \mathrm{IX}, 1, 7\).

\(\mu^*(f), \mu^*(\mathrm{A}), \int^* f d\mu, \int^* f(t) d\mu(t): \mathrm{IX}, 1, 9\).

\(\mu^*: \mathrm{IX}, 1, 9\).

\(\overline{\mathcal{L}}^p(\mathrm{T}, \mu), \overline{\mathcal{L}}_F^p(\mathrm{T}, \mu), \mathcal{L}^p(\mathrm{T}, \mu), \mathcal{L}_F^p(\mathrm{T}, \mu)\) (for \(1 \leqslant p \leqslant +\infty\)): IX, 1, 10.

\(\overline{\mathcal{L}}_F^p(\mu), \overline{\mathcal{L}}_F^p, \overline{\mathcal{L}}^p, \overline{\mathcal{L}}^p(\mu), \mathcal{L}^p(\mu), \mathcal{L}^p: \mathrm{IX}, 1, 10\).

\(\overline{\mathrm{N}}_p(f), \mathrm{N}_p(f), \overline{\mathcal{N}}_F, \mathcal{N}_F: \mathrm{IX}, 1, 10\).

\(\mathrm{L}_F^p(\mu), \mathrm{L}_F^p: \mathrm{IX}, 1, 10\).

\(\int f d\mu, \mu(f), \int f(t) d\mu(t): \mathrm{IX}, 1, 10\).

\(\mu_X^\bullet, \mu_X, \mu|_{\mathrm{X}}: \mathrm{IX}, 2, 1\).

\(f \cdot \mu: \mathrm{IX}, 2, 2\).

\(\pi(\mu): \mathrm{IX}, 2, 3\).

\(\lambda \otimes \mu: \mathrm{IX}, 2, 5\).

\(\Re(\mathrm{T}), \Re(\mathrm{T}): \text{conventions of §3}\).

\(\mathcal{C}^b(\mathrm{T}; \mathrm{F}), \mathcal{C}^b(\mathrm{T}), \mathcal{C}^b, \mathcal{C}_+^b(\mathrm{T}), \mathcal{C}_+^b: \text{conventions of §5}\).

\(\mathcal{M}^b(\mathrm{T}; \mathbf{C}), \mathcal{M}^b(\mathrm{T}), \mathcal{M}^b, \mathcal{M}_+^b(\mathrm{T}), \mathcal{M}_+^b: \text{conventions of §5}\).

\(\mathcal{L}\mu: \mathrm{IX}, 5, 7\).

\(\mathcal{F}(\mathrm{E}): \mathrm{IX}, 6, 1\).

\(p_V, p_{VW}: \mathrm{IX}, 6, 1\).

\(\mathcal{D}(\mathrm{E}): \mathrm{IX}, 6, 1\).

\(\widetilde{\lambda}: \mathrm{IX}, 6, 1\).

\(u(\mu)\) (\(\mu\) a promeasure): IX, 6, 2.

\(\mathcal{F}\mu\) (\(\mu\) a promeasure or a measure): IX, 6, 3.

\(\Gamma_Q, \gamma_a: \mathrm{IX}, 6, 5\).

\(\gamma_C: \mathrm{IX}, 6, 6\).

\(\operatorname{Tr}(\mathrm{Q}/\mathrm{H}): \mathrm{IX}, \text{Annex}, 1\).

\(u^*: \mathrm{IX}, \text{Annex}, 2\).
